\chapter{ProofTools — Representatives, extensions and comparison laws}
Fix a model $D$ of the classical and synthetic Adams systems. All maps, quotient objects, filtrations and page labels in each statement are taken in this model. In the statements quoted from \cite{LWX126}, spectra are $2$-completed, connective and of finite type.

For the level-indexed rules, page-extension comparisons and the
composition theorem, $D$ is the fixed Standard model and the classical
objects belong to its selected family. Statements about general filtered
groups or convergent towers specify their own data explicitly.

The page-extension language is fixed in
Definitions~\ref{paper:5.4} and~\ref{sem:AdamsSynthetic:page-extension-levels}.
We compare zero-level and positive-level relations with their actual maps,
weights and target indeterminacies. Examples of page extensions and
crossings accompany their definitions.
Classical realization is discussed with its target-preserving lifting
criterion. The Leibniz rule is followed by the warning about crossing,
and the Mahowald rule by its cofiber examples. Restriction and stretching
precede the original infinite-extension application and its Leibniz
calculation. The general finite-quotient composition is then developed
independently of these numerical applications. The final section treats
Adams differential transport and the classical realization of finite-level
extensions, with their actual crossing and permanence hypotheses.

\section{Common labels and representative calculus}

\begin{definition}[Common labels and differential equations]
\label{sem:ProofTools:common-labels}
\level{coarse}
For a classical spectrum $X$, a common label in bidegree $(s,t)$ is an element of $E_2^{s,t}(X)$. For $c\geq1$, it is a cycle through $d_c$, respectively a boundary through $d_c$, if it lies in $Z_c^{s,t}(X)$, respectively $B_c^{s,t}(X)$. A permanent label is an element of $Z_\infty^{s,t}(X)$; the zero label is included. The index $c$ is distinct from the lambda level of a synthetic label.

For $c\geq1$, labels $x,y\in E_2^{s,t}(X)$ are congruent modulo $B_c$ if $x-y\in B_c^{s,t}(X)$. For $r\geq2$, the equation $d_r(x)=y$ means that $x$ and $y$ represent classes on page $r$ and their classes satisfy the Adams differential equation, with source bidegree $(s,t)$ and target bidegree $(s+r,t+r-1)$. The target may be zero.
\lean{KIP126.ProofTools.Ambient,KIP126.ProofTools.IsCycle,KIP126.ProofTools.IsBoundary,KIP126.ProofTools.IsPermanent,KIP126.ProofTools.Congruent,KIP126.ProofTools.DifferentialAt}
\end{definition}

For $r\geq2$, a common label lies in $Z_{r-1}$ if and only if it reaches page $r$. If $x$ represents $x_r$ on page $r$ and $y-x\in B_{r-1}$, then $y$ represents $x_r$. If $d_r(x)=y$ and $y'-y\in B_{r-1}$, then $d_r(x)=y'$.

For $c\geq1$, $B_c\subseteq Z_c$, and $B_1=0$. If a common label represents zero on page $r\geq2$, it lies in $B_{r-1}$. Projection of a common cycle to the initial page recovers its $E_2$-label.

\begin{theorem}[Filtration, associated grades and detection]
\label{sem:ProofTools:filtration-representatives}
\level{coarse}
For a convergent synthetic Adams tower on $X$, let
\[E_\infty^{s,m+s,w}(X)\cong\operatorname{gr}^s\pi_{m,w}X\]
be its convergence isomorphism.

Every element of $F^s\pi_{m,w}X$ has a representative in the filtration subobject $F^sX$. If $g:X\to Y$ is a map of convergent towers, then
\[g_*\bigl(F^s\pi_{m,w}X\bigr)\subseteq F^s\pi_{m,w}Y.\]

Let $g:X\to Y$ be a map of convergent towers, $a\in F^s\pi_{m,w}X$ and $b\in F^s\pi_{m,w}Y$. If $g_*(a)=b$, their associated-graded classes satisfy $g_\infty([a])=[b]$. Conversely, this graded equality implies
\[b-g_*(a)\in F^{s+1}\pi_{m,w}Y.\]
If $g_\infty$ is injective in this degree and $g_*(a)\in F^{s+1}\pi_{m,w}Y$, then $a\in F^{s+1}\pi_{m,w}X$.

Let $g:X\to Y$ be a map of convergent towers. If $x$ detects $\alpha$ on $X$, then $g_2(x)$ detects $g_*(\alpha)$ on $Y$. Conversely, suppose that $g_\infty$ is injective in filtration $s$, that $x$ represents an infinite-page class, and that $\alpha\in F^s\pi_{m,w}X$. If $g_2(x)$ detects $g_*(\alpha)$, then $x$ detects $\alpha$.
\uses{sem:AdamsSynthetic:selected-convergence,sem:SpectralSequence:filtration,sem:SpectralSequence:detection,sem:AdamsSynthetic:families}
\lean{KIP126.ProofTools.Internal.filtration_representative_exists,KIP126.ProofTools.Internal.map_preserves_filtration,KIP126.ProofTools.Internal.grade_natural,KIP126.ProofTools.Internal.grade_difference,KIP126.ProofTools.Internal.next_filtration_of_infinity_injective,KIP126.ProofTools.Internal.detects_map,KIP126.ProofTools.Internal.detects_reflect}
\end{theorem}
\begin{proof}
Filtration classes are images of the corresponding tower stages, so they have representatives there and a map of towers preserves them. The kernel of passage from $F^s$ to its associated grade is $F^{s+1}$. Applying this observation to $b-g_*(a)$ proves naturality and the next-filtration ambiguity; injectivity on the associated grade gives the stated reflection.

Under the convergence isomorphisms, detection is precisely equality of associated-graded classes. Naturality transports that equality, and the assumed injectivity reflects it. The induced page maps compose by functoriality, and send an infinite-page representative to a representative of its image. These are statements about the specified convergent towers.
\end{proof}

\begin{theorem}[Representatives for a tower extension equation]
\label{sem:ProofTools:extension-representatives}
\level{coarse}
Let $g:X\to Y$ be a map of convergent towers. Fix $(m,w)$ and classes $x\in E_\infty^{s,m+s,w}(X)$ and $y\in E_\infty^{s+n,m+s+n,w}(Y)$. An extension equation of length $n$ from $x$ to $y$ holds if and only if $n\geq0$ and there are representatives
\[a\in F^s\pi_{m,w}X,\qquad b\in F^{s+n}\pi_{m,w}Y,\qquad g_*(a)=b\]
whose associated-graded classes correspond to $x$ and $y$ under the convergence isomorphisms.
\uses{sem:ProofTools:filtration-representatives,paper:2.3}
\lean{KIP126.ProofTools.Internal.tower_extension_iff}
% Blueprint note: This is the representative criterion for the specified towers. Identifying an arbitrary classical extension spectral sequence with this tower construction remains a separate comparison; this entry does not assert that comparison.
\end{theorem}
\begin{proof}
Unpack the extension relation in the filtered homotopy groups of the two towers. Its source and target grades are represented by $a$ and $g_*(a)$ in the displayed filtrations; conversely such representatives give the extension equation. The convergence identifications translate these grades into the specified infinite-page classes.
\end{proof}

\section{Finite filtration correction and quotient windows}

\begin{theorem}[Finite correction of a residual]
\label{sem:ProofTools:finite-correction}
\level{coarse}
Let $A,A',B$ be abelian groups, $N\in\mathbb N$, and $\rho:A'\to A$ and $g:A\to B$ homomorphisms, let $E_k\leq A$, $L_k\leq A'$ with $L$ decreasing, and let $T\leq B$. Suppose that for each $s\leq k<s+N$, a residual $\delta\in E_k$ can either be corrected by some $\varepsilon\in L_k$ so that $\delta-\rho\varepsilon\in E_{k+1}$, or already satisfies $g\delta\in T$. Suppose also that $g(E_{s+N})\subseteq T$. If $\alpha_0-\rho\alpha\in E_s$, there is $\alpha'$ with
\[\alpha'-\alpha\in L_s,\qquad g(\alpha_0-\rho\alpha')\in T.\]
\lean{KIP126.ProofTools.finite_filtration_correction}
% Blueprint note: The argument uses only a finite interval; it constructs neither a completion nor a compatible infinite lift.
\end{theorem}
\begin{proof}
Induct over the finite interval. If the current residual already has image in $T$, keep the current correction. Otherwise subtract the supplied $\rho\varepsilon$ and add $\varepsilon$ to the source representative, thereby moving the residual to the next filtration. Since $L$ decreases, every accumulated change lies in $L_s$. At the terminal stage the hypothesis on $g(E_{s+N})$ puts its image in $T$. This proves both conclusions without an infinite limiting process.
\end{proof}

\begin{theorem}[Finite quotient windows and restriction]
\label{sem:ProofTools:finite-strip}
\level{coarse}
For $q>0$, the infinite page of $\nu X/\lambda^q$ vanishes outside $0\leq t-w<q$. Every $\alpha\in\pi_{t-s,t}(\nu X/\lambda^q)$ has filtration at least $s$. If $q\leq m+s-w$, an element of $\pi_{m,w}(\nu X/\lambda^q)$ in filtration at least $s$ is zero.

For $0<i\leq j$ and $t-w<i$, the quotient reduction $\rho:\nu X/\lambda^j\to\nu X/\lambda^i$ induces an injective map on the infinite-page component in tridegree $(s,t,w)$.
\uses{sem:Literature:infinite-finite-window,sem:Literature:infinite-restriction,sem:AdamsSynthetic:finite-quotient-separation,sem:ProofTools:filtration-representatives}
\lean{KIP126.ProofTools.Internal.Mahowald.finite_infinity_zero_outside,KIP126.ProofTools.Internal.Mahowald.finite_bottom_filtration,KIP126.ProofTools.Internal.Mahowald.finite_high_filtration_zero,KIP126.ProofTools.Internal.rho_infty_injective}
\end{theorem}
\begin{proof}
The finite-window comparison gives the stated support of the infinite page. Vanishing associated grades identify successive filtration subgroups; iterating over a finite interval gives the corresponding filtration bounds. Finite-quotient separation then gives the asserted zero subgroup above the upper edge of the strip. This separation is used only for the finite quotient.

In the retained strip, the two quotient windows have the same boundary denominator, and reduction is induced by inclusion of the classical cycle groups. It is therefore injective. Equivalently, filtration reflection can be checked one associated grade at a time over the finite interval; no injectivity of the whole homotopy map is asserted.
\end{proof}

\begin{theorem}[Finite target labels and sphere detection]
\label{sem:ProofTools:sphere-label-detection}
\level{coarse}
For $q>0$ and $0\leq k<q$, let $y\in Z_{q-k}^{s,t}(X)$. Its finite target label at weight $t-k$ represents $e\in E_\infty^{s,t,t-k}(\nu X/\lambda^q)$ if and only if the finite-window comparison carries $e$ to the class of $y$ in
\[Z_{q-k}^{s,t}(X)/B_{1+k}^{s,t}(X).\]

The first-quotient comparison carries the first label of $y\in E_2^{s,t}(S^0)$ to $y$. Let $\alpha\in\pi_{t-s,t}(\nu S^0)$ have first-quotient image equal to this label, using the canonical zero-shift identification. For $q>0$ and $k\geq0$, the finite target label of $y$ in tridegree $(s,t,t-k)$ detects the image of $\lambda^k\alpha$ in $\nu S^0/\lambda^q$.
\uses{sem:ProofTools:filtration-representatives,sem:ProofTools:finite-strip,sem:Literature:infinite-finite-window,sem:Literature:infinite-shift,sem:AdamsSynthetic:e2-comparison,sem:AdamsSynthetic:first-quotient-shifts,sem:AdamsSynthetic:ess-shifts}
\lean{KIP126.ProofTools.Internal.Mahowald.finite_target_label_iff,KIP126.ProofTools.Internal.Mahowald.sphere_first_quotient_label,KIP126.ProofTools.Internal.Mahowald.sphere_finite_lambda_detects}
\end{theorem}
\begin{proof}
The finite-window isomorphism identifies a specified finite target label with precisely the displayed classical cycle coset. At the first quotient the fixed sphere comparison recovers the original label. Naturality of quotient projection and multiplication by $\lambda$ transports this detection to each indicated weight and finite quotient. The statement retains the specified label and does not impose nonzero detection on its image.
\end{proof}

\section{Actual crossings and common representatives}

\begin{definition}[Actual essential no-crossing]
\label{sem:ProofTools:mh-actual-no-crossing}
\level{coarse}
Fix a map $F:U\to V$ of convergent synthetic Adams towers and one stem and weight $(p,w)$. For a source filtration $s$ and length $a$, actual no-crossing means
\[
\text{there is no essential }d_{v-u}^{F}(\xi)=\eta
\quad\text{with }s<u\leq v\leq s+a
\]
in this same map and bidegree. Essentiality is taken in the actual extension spectral sequence, including all shorter extension images. The diagonal $u=v$ is included.
\uses{sem:AdamsSynthetic:selected-convergence,sem:SpectralSequence:filtration}
\lean{KIP126.ProofTools.Internal.Mahowald.ActualNoCrossing}
\end{definition}

\begin{theorem}[A common source from actual no-crossing]
\label{sem:ProofTools:mh-common-source}
\level{coarse}
Let $F:U\to V$ and $G:U\to W$ be maps of convergent towers with the same source convergence. Fix the same stem and weight, and suppose
\[
d_a^F(\xi)=\eta,\qquad d_b^G(\xi)=\zeta.
\]
If either equation has actual no-crossing in its full higher-source triangle, then there are actual representatives $u,v,w$ of $\xi,\eta,\zeta$, respectively, with
\[F(u)=v,\qquad G(u)=w.\]
The representatives lie in filtrations $s,s+a,s+b$. Zero labels, inessential input equations, and length zero are allowed.

The shared representative is also the issue emphasized by \cite{LWX126}, Remark 6.18, in its discussion of a Mahowald rule without crossing hypotheses. Two separately witnessed input relations do not specify a common representative with both prescribed images. The conclusion here produces that common representative under the stated actual no-crossing alternative; it does not assert a rule with that hypothesis omitted.
\uses{sem:ProofTools:mh-actual-no-crossing,sem:ProofTools:filtration-representatives,sem:ProofTools:finite-correction}
\lean{KIP126.ProofTools.Internal.Mahowald.actual_common_source_representative}
\end{theorem}
\begin{proof}
\textit{Proof sketch in the fixed model.}
Choose representatives for the two given extension equations. Their source representatives differ in the next source filtration. Synchronize them over the finite higher-source triangle: a first nonzero image grade of a residual that cannot be removed gives an essential crossing of the edge on which actual no-crossing is assumed. The finite correction argument therefore supplies a common source while retaining the two prescribed image grades. The same argument includes diagonal extensions and zero labels; essentiality of the input equations is not used.
\end{proof}

Assume that suspension and bigraded shifts are coherent with the homotopy functors. The canonical isomorphism
\[\Sigma^{1,-q}Q\cong\Sigma(\Sigma^{0,-q}Q)\]
identifies the induced suspension and bigraded-shift maps. The inverse and composition isomorphisms for shifts induce the corresponding inverse and composition maps on homotopy groups.
For all integers $m,n,k,l$ and $f,g\in\pi_{m,n}X$, suspension invariance is additive:
\[
\operatorname{susp}_{k,l}(f+g)=\operatorname{susp}_{k,l}(f)+\operatorname{susp}_{k,l}(g)
\quad\text{in }\pi_{m+k,n+l}(\Sigma^{k,l}X).
\]

\section{No-crossing criteria for filtered extensions}

% AimPaper: 2.6 corollary 327 330
\begin{corollary}[Vanishing below the Adams filtration]
\label{paper:2.6}
\source{AimPaper Corollary 2.6}
\level{coarse}
From \cite{LWX126}, Corollary 2.6.\par
% BEGIN AIMPAPER STATEMENT
If the Adams filtration of $f$ is $k$, then
    $d^f_i=0$ for $i<k$.
% END AIMPAPER STATEMENT
\uses{paper:2.5,sem:AdamsSynthetic:filtration-comparison}
% Blueprint note: Mathematical dependency review: [AimPaper source, lines 332-339](paper/main.tex.html#L332). Factor f into positive-filtration maps, then apply Proposition 2.5(1) to rule out shorter extension differentials.

% Correspondence: fixed original representative-form statement proved.
% Blueprint note: Original statement and context: paper/main.tex.html#L327; paper/2412.10879.pdf.
% Blueprint note: Original-statement interface: \texttt{KIP126.ProofTools.Paper.adamsFiltrationVanishing}. The public statement is proved without changing its type or premises.
% Blueprint note: The proof lifts the map along the actual target Adams tower and proves that filtration adds under composition. Every shorter extension target has zero associated grade.
\lean{KIP126.ProofTools.Paper.adamsFiltrationVanishing}
\end{corollary}
\begin{proof}
\textit{(See \cite{LWX126}, Corollary 2.6.)} \href{paper/main.tex.html#L332}{Original proof}.

When $k>0$, it is clear that $d^f_0=0$, which means that for any $s$ and $[x]\in F_s\pi_*(X)$, we have $f([x])\in F_{s+1}\pi_*(Y)$. Since $\AF(f)=k$, we know that there exist $k$ maps
    $$X=X_0\fto{f_1}X_1\fto{f_2}\cdots\fto{f_k}X_k=Y$$
    such that each $f_i$ induces trivial map in $\HF$-homology (hence $\AF(f_i)>0$) and $f$ is homotopic to the composition $f_k\circ \cdots\circ f_1.$
    Then for any $s$ and $[x]\in F_s\pi_*(X)$, inductively we have
    $$f_1([x])\in F_{s+1}\pi_*(X_1)$$
    $$f_2\circ f_1([x])\in F_{s+2}\pi_*(X_2)$$
    $$\cdots$$
    $$f([x])=f_k\circ\cdots\circ f_1([x])\in F_{s+k}\pi_*(X_k)=F_{s+k}\pi_*(Y).$$
    By Proposition \ref{prop:8154e6f1}.(1) this means that $d_i^fx=0$ for any $x\in E_\infty^{*,*}(X)$ and $i<k$.
\end{proof}

% AimPaper: 2.7 definition 344 352
\begin{definition}[Crossings of extensions]
\label{paper:2.7}
\source{AimPaper Definition 2.7}
\level{coarse}
From \cite{LWX126}, Definition 2.7.\par
% BEGIN AIMPAPER STATEMENT
\label{def:98skj23}
\begin{enumerate}[leftmargin=*]
    \item For $x\in E_\infty^{s,t}(X)$, we say that the an $f$-extension $d_n^{f}(x)=y$ has a crossing that hits Adams filtration $p$ for some $p \leq s+n$, if there exists an element $x^\prime\in E_\infty^{s+a,t+a}(X)$ with $a>0$ and an $f$-extension $d_m^{f}(x^\prime)=y^\prime$ for $0\neq y^\prime\in E^{s+a+m, t+a+m}_\infty(Y)$ such that
    $$p = \textup{AF}(y^\prime) = s+a+m\le \textup{AF}(y) = s+n.$$
    \item We say that an $f$-extension $d_n^{f}(x)=y$ has no crossing that hits the range AF $\ge p$ if there does not exist an element $x^\prime\in E_\infty^{s+a,t+a}(X)$ with $a>0$ and an $f$-extension $d_m^{f}(x^\prime)=y^\prime$ for $0\neq y^\prime\in E^{s+a+m, t+a+m}_\infty(Y)$ such that
    $$p\le \textup{AF}(y^\prime) = s+a+m\le \textup{AF}(y) = s+n.$$
    \item We say that an $f$-extension $d_n^{f}(x)=y$ has no crossing if it has no crossing that hits the range AF$\ge \AF(x) + 1$.
\end{enumerate}
% END AIMPAPER STATEMENT

Let $f$ be a map of convergent towers or filtered abelian groups. An extension from source filtration $s$ to target filtration $s+n$ has no higher-source crossing over $[u,s+n]$ if every extension with source filtration greater than $s$ and target filtration in $[u,s+n]$ has zero target label. All extensions, including inessential ones, occur in this quantification.
\uses{paper:2.3}
% Blueprint note: Mathematical dependency review: [AimPaper source, lines 344-352](paper/main.tex.html#L344). A crossing is a competing extension from a higher-filtration source, using the extension relation already defined.

% Blueprint note: The generic filtered-map crossing predicate belongs to \ref{sem:SpectralSequence:filtered-transport}. The bindings here are the two stated no-crossing predicates with an arbitrary filtration threshold.

% Correspondence: partial.
% Blueprint note: Public representative-level definitions \texttt{FilteredNoCrossingRange} and \texttt{ExtensionNoCrossingRange} allow arbitrary filtration thresholds. Their comparison with \texttt{DifferentialDatum} and \texttt{HasCrossingAt} for a general classical $f$-extension spectral sequence has not yet been assembled into a complete public bridge, so this correspondence remains partial.
% Blueprint note: Direct formal counterparts: \texttt{KIP126.ProofTools.Paper.FilteredNoCrossingRange}, \texttt{KIP126.ProofTools.Paper.ExtensionNoCrossingRange}.
\lean{KIP126.ProofTools.Paper.FilteredNoCrossingRange,KIP126.ProofTools.Paper.ExtensionNoCrossingRange}

\sketch
% Blueprint note: Original statement and context: paper/main.tex.html#L344; paper/2412.10879.pdf.
\end{definition}

\begin{theorem}[The tower representative criterion]
\label{sem:ProofTools:representative-criterion}
\level{coarse}
From \cite{LWX126}, Proposition 2.10.\par

Fix a map of convergent towers and an extension equation of length $n$ from $x$ to $y$ in filtrations $s$ and $s+n$, and a threshold $u\leq s+n$. Absence of higher-source extensions whose target lies between $u$ and $s+n$ is equivalent to the following condition: every representative $a$ of $x$ whose image has filtration at least $u$ has image representing the specified $y$ in filtration $s+n$. For threshold $u=s$, the image-filtration premise is automatic. If $y=0$, the conclusion says that every such image has filtration at least $s+n+1$. A map increasing every filtration by at least $n$ satisfies this criterion.
\uses{sem:ProofTools:extension-representatives,paper:2.7,sem:ProofTools:finite-correction}
\lean{KIP126.ProofTools.Paper.extension_noCrossingRange_iff,KIP126.ProofTools.Paper.extension_noCrossing_iff,KIP126.ProofTools.Paper.extension_zero_noCrossing_iff,KIP126.ProofTools.Paper.extension_noCrossing_of_filtration_bound}
% Blueprint note: The bindings here are the four tower representative forms. The filtered-group forms remain at Proposition 2.10. No new comparison with an arbitrary classical extension spectral sequence is asserted.
\end{theorem}
\begin{proof}
Use representatives as in the proof of Proposition~\ref{paper:2.10}. If two representatives of the source have different image grades in the specified range, their difference is in higher source filtration and produces a crossing; a first nonzero grade is obtained over the finite interval. Conversely, adding a representative of a crossing source preserves the original source grade but changes the required target grade. This proves the equivalence. The zero-target case and the filtration-bound case follow by reading the target grade as a zero associated-graded class.
\end{proof}

% AimPaper: 2.10 proposition 408 416
\begin{proposition}[Characterization of extensions without crossings]
\label{paper:2.10}
\source{AimPaper Proposition 2.10}
\level{coarse}
From \cite{LWX126}, Proposition 2.10.\par
% BEGIN AIMPAPER STATEMENT
\label{prop:i8r47oe}
    An $f$-extension from $x$ to $y$ has no crossing in the range $\AF\ge p$ if and only if for all $[x]\in \{x\}$ such that $\AF(f[x])\ge p$ we have $f[x]\in \{y\}$. In particular, we have
    \begin{enumerate}
        \item An $f$-extension from $x$ to $y$ has no crossing if and only if for all $[x]\in \{x\}$ we have $f[x]\in \{y\}$.
        \item An $f$-extension $d_n^{f}(x)=0$ has no crossing if and only if for all $[x]\in \{x\}$,
        $$\AF(f[x]) >\AF(x)+n.$$
        \item If $\AF(f)=n$, then all $d^f_n$-differentials have no crossing in the $f$-ESS.
    \end{enumerate}
% END AIMPAPER STATEMENT
\uses{paper:2.5,paper:2.7}
% Blueprint note: Mathematical dependency review: [AimPaper source, lines 418-430](paper/main.tex.html#L418). Proposition 2.5(3) and the definition of a crossing turn universal representative behavior into the no-crossing criterion.

% Correspondence: public counterpart, within the scope described below.
% Blueprint note: The full representative-level equivalence and its three clauses are public for arbitrary filtered abelian groups. An $E_\infty$ formulation for the same synthetic map is also available. This does not establish the missing general classical $f$-extension spectral sequence construction in Definition~2.1.
% Blueprint note: Formal counterparts: \item \texttt{KIP126.ProofTools.Paper.filtered\_extension\_noCrossingRange\_iff} \item \texttt{KIP126.ProofTools.Paper.filtered\_extension\_noCrossing\_iff} \item \texttt{KIP126.ProofTools.Paper.filtered\_zero\_extension\_noCrossing\_iff} \item \texttt{KIP126.ProofTools.Paper.filtered\_extension\_noCrossing\_of\_filtration\_bound}
\lean{KIP126.ProofTools.Paper.filtered_extension_noCrossingRange_iff,KIP126.ProofTools.Paper.filtered_extension_noCrossing_iff,KIP126.ProofTools.Paper.filtered_zero_extension_noCrossing_iff,KIP126.ProofTools.Paper.filtered_extension_noCrossing_of_filtration_bound}

% Blueprint note: Original statement and context: paper/main.tex.html#L408; paper/2412.10879.pdf.
\end{proposition}
\begin{proof}
\textit{(See \cite{LWX126}, Proposition 2.10.)} \href{paper/main.tex.html#L418}{Original proof}.

First we prove the only if part. Assume by contradiction that there exists $[x]\in \{x\}$ such that $f[x]=[y']\notin \{y\}$ is detected by $y'\in E_\infty(Y)$ with $\AF(y')\ge p$. Then there is an $f$ extension from $x$ to $y'$. By Proposition \ref{prop:8154e6f1}~(3), $y'$ (if $\AF(y)>\AF(y')$)  or $y-y'$ (if $\AF(y)=\AF(y')$) or $y$ (if $\AF(y)<\AF(y')$) should be hit by a shorter $d^f$ differential. This shorter differential is a crossing of the $f$-extension from $x$ to $y$ that hits $\AF\ge p$, which is a contradiction. Now we prove the if part. Assume that the $f$-extension has a crossing from $x'$ to $y'$ with $p\le \AF(y')\le \AF(y)$. Then there exists $[x']\in \{x'\}$ such that $f([x'])\in \{y'\}$. Note that $[x]+[x']\in \{x\}$ since $\AF(x')>\AF(x)$, we have
    $$f([x]+[x'])=f([x])+f([x'])=[y]+[y']\in\begin{cases}
        \{y'\} & \text{ if } \AF(y')<\AF(y)\\
        \{y+y'\} & \text{ if } \AF(y')=\AF(y).\\
    \end{cases}$$
    However, by assumption we should have $f([x]+[x'])\in \{y\}$ and it is contradictory to the equation above in either case.

    Hence the main statement is true.  Part (1) (2) are special cases of the main statement and Part (3) holds for degree reasons.
\end{proof}

\section{Commutative squares and extension comparison}

% AimPaper: 2.12 theorem 461 481
\begin{theorem}[The extension theorem for a commutative square]
\label{paper:2.12}
\source{AimPaper Theorem 2.12}
\level{coarse}
From \cite{LWX126}, Theorem 2.12.\par
% BEGIN AIMPAPER STATEMENT
\label{thm:4114f70c}
    Consider a homotopy commutative diagram of spectra
    $$\xymatrix{
    X \ar[r]^{f}\ar[d]_{p} & Y\ar[d]^{q}\\
    Z \ar[r]_{g} & W
    }$$
    Suppose $m,n,l\ge 0$, $0<k\le m+l-n$,
    $$\begin{array}{ll}
        x\in E^{s,t}_\infty(X) & y\in E^{s+n,t+n}_\infty(Y),\\
        z\in E^{s+m,t+m}_\infty(Z) & w\in E^{s+m+l,t+m+l}_\infty(W),
    \end{array}$$
    and
    \begin{enumerate}
        \item $d^{f}_{n}(x)=y$,
        \item $d^{p}_{m}(x)=z$,
        \item the differential in (1) or the differential in (2) has no crossing,
        \item $d^{g}_{l}(z)=w$, and has no crossing that hits the range AF $\ge s+n+k$,
        \item $d^{q}_{k-1} y=0$, and has no crossing.
    \end{enumerate}
    Then $d^{q}_{m+l-n}(y)=w$. (See Figure \ref{fig:4f154a8}.)
% END AIMPAPER STATEMENT
\uses{paper:2.10}
% Blueprint note: Mathematical dependency review: [AimPaper source, lines 483-505](paper/main.tex.html#L483). Use the representative criterion to choose a common source and transport its images around the commutative square.

% Correspondence: fixed original representative-form statement proved.
% Blueprint note: Original statement and context: paper/main.tex.html#L461; paper/2412.10879.pdf.
% Blueprint note: Original-statement interface: \texttt{KIP126.ProofTools.Paper.extensionSquare}. The public statement is proved without changing its type or premises.
% Blueprint note: The proved finite-filtration representative criterion synchronizes the two source edges. The zero $q$-extension puts the same representative above $s+n+k$; only the stated threshold no-crossing condition for $g$ is then used.
\lean{KIP126.ProofTools.Paper.extensionSquare}
\end{theorem}
\begin{proof}
\textit{(See \cite{LWX126}, Theorem 2.12.)} \href{paper/main.tex.html#L483}{Original proof}.

First we find a representative $[x]$ of $x$ such that
    \begin{equation}\label{eq:9d45feac}
        p[x]\in \{z\}
    \end{equation}
    and
    \begin{equation}\label{eq:36e72616}
        f[x]\in \{y\}.
    \end{equation}
    If (1) has no crossing, by (2) we can pick $[x]\in \{x\}$ such that (\ref{eq:9d45feac}) holds;
    If (2) has no crossing, by (1) we can pick $[x]\in \{x\}$ such that  (\ref{eq:36e72616}) holds.
    In both cases the other equation also holds because of the no crossing condition.

    By (5), (\ref{eq:36e72616}) and Proposition~\ref{prop:i8r47oe}~(2) we know that $qf[x]$ has AF $\ge s+n+k$. Since the diagram on the Adams $E_\infty$-pages commutes, we have
    $$AF(gp[x])=AF(qf[x])\ge s+n+k.$$
    Combining with (4), (\ref{eq:9d45feac}) and Proposition~\ref{prop:i8r47oe} we have
    $$gp[x]\in \{w\}.$$
    Therefore, $qf[x]\in \{w\}$ and by (\ref{eq:36e72616}) there is a $q$-extension from $y$ to $w$.
\end{proof}

% AimPaper: 2.15 corollary 555 574
\begin{corollary}[A square of extensions without crossings]
\label{paper:2.15}
\source{AimPaper Corollary 2.15}
\level{coarse}
From \cite{LWX126}, Corollary 2.15.\par
% BEGIN AIMPAPER STATEMENT
\label{cor:0012nik}
    Consider a homotopy commutative diagram of spectra
    $$\xymatrix{
    X \ar[r]^{f}\ar[d]_{p} & Y\ar[d]^{q}\\
    Z \ar[r]_{g} & W
    }$$
    Suppose $m,n,l\ge 0$, 
    $$\begin{array}{ll}
        x\in E^{s,t}_\infty(X) & y\in E^{s+n,t+n}_\infty(Y),\\
        z\in E^{s+m,t+m}_\infty(Z) & w\in E^{s+m+l,t+m+l}_\infty(W),
    \end{array}$$
    and
    \begin{enumerate}
        \item $d^{f}_{n}(x)=y$,
        \item $d^{p}_{m}(x)=z$,
        \item the differential in (1) or the differential in (2) has no crossing,
        \item $d^{g}_{l}(z)=w$, and has no crossing.
    \end{enumerate}
    Then $d^{q}_{m+l-n}(y)=w$.
% END AIMPAPER STATEMENT
\uses{paper:2.12}
% Blueprint note: Mathematical dependency review: [AimPaper source, lines 555-574](paper/main.tex.html#L555). Specialize the square theorem to the full no-crossing hypothesis on the lower extension.

% Correspondence: fixed original representative-form statement proved.
% Blueprint note: Original statement and context: paper/main.tex.html#L555; paper/2412.10879.pdf.
% Blueprint note: Original-statement interface: \texttt{KIP126.ProofTools.Paper.noCrossingExtensionSquare}. The public statement is proved without changing its type or premises.
% Blueprint note: The proved filtered-group no-crossing criterion is specialized directly to the actual classical Adams tower. It synchronizes a common source representative and the lower target, and the same commutative square gives the right extension. No synthetic-to-classical comparison or additional premise is needed.
\lean{KIP126.ProofTools.Paper.noCrossingExtensionSquare}
\end{corollary}
\begin{proof}
\textit{(See \cite{LWX126}, Corollary 2.15.)} \href{paper/main.tex.html#L555}{Original statement and context}.

Choose one representative $[x]$ for which both $f[x]$ and $p[x]$ have the prescribed detectors $y$ and $z$. Such a common choice exists because one of these two extensions has no crossing: choose a representative for the other extension, and apply Proposition~\ref{paper:2.10} to the crossing-free one.

The crossing-free $g$-extension then implies that $gp[x]$ is detected by $w$. Homotopy commutativity gives $qf[x]=gp[x]$. Thus $f[x]$, detected by $y$, has $q$-image detected by $w$, which is exactly the required extension. Its length is the filtration difference $(s+m+l)-(s+n)=m+l-n$.
\end{proof}

% AimPaper: 2.16 corollary 576 591
\begin{corollary}[Extensions in a commutative triangle]
\label{paper:2.16}
\source{AimPaper Corollary 2.16}
\level{coarse}
From \cite{LWX126}, Corollary 2.16.\par
% BEGIN AIMPAPER STATEMENT
\label{cor:166dc180}
    Consider a homotopy commutative diagram of spectra
    $$\xymatrix{
    X \ar[r]^{f}\ar[dr]_{p} & Y\ar[d]^{q}\\
    & Z
    }$$
    Suppose $n,m\ge 0$, 
    $$x\in E^{s,t}_\infty(X), ~y\in E^{s+n,t+n}_\infty(Y), ~z\in E^{s+m,t+m}_\infty(Z)$$
    and
    \begin{enumerate}
        \item $d^{f}_{n}(x)=y$,
        \item $d^{p}_{m}(x)=z$,
        \item the differential in (1) or the differential in (2) has no crossing.
    \end{enumerate}
    Then $d^{q}_{m-n}(y)=z$.
% END AIMPAPER STATEMENT
\uses{paper:2.15}
% Blueprint note: Mathematical dependency review: [AimPaper source, lines 593-595](paper/main.tex.html#L593). Specialize Corollary 2.15 to an identity lower horizontal map.

% Correspondence: fixed original representative-form statement proved.
% Blueprint note: Original statement and context: paper/main.tex.html#L576; paper/2412.10879.pdf.
% Blueprint note: Original-statement interface: \texttt{KIP126.ProofTools.Paper.extensionTriangle}. The public statement is proved without changing its type or premises.
% Blueprint note: The common-source representative criterion proves the third extension directly. The conclusion keeps length $m-n$, including the signed-length cases of the fixed statement.
\lean{KIP126.ProofTools.Paper.extensionTriangle}
\end{corollary}
\begin{proof}
\textit{(See \cite{LWX126}, Corollary 2.16.)} \href{paper/main.tex.html#L593}{Original proof}.

This is a special case of Corollary \ref{cor:0012nik} when $Z=W$ and $g=id_Z$.
\end{proof}

% AimPaper: 2.17 corollary 597 611
\begin{corollary}[Composition of extensions]
\label{paper:2.17}
\source{AimPaper Corollary 2.17}
\level{coarse}
From \cite{LWX126}, Corollary 2.17.\par
% BEGIN AIMPAPER STATEMENT
\label{cor:290d35ce}
    Consider a homotopy commutative diagram of spectra
    $$\xymatrix{
    X \ar[d]_{p}\ar[dr]^{q} & \\
    Z\ar[r]_{g} & W
    }$$
    Suppose $n,m\ge 0$, 
    $$x\in E^{s,t}_\infty(X), ~z\in E^{s+m,t+m}_\infty(Y), ~w\in E^{s+m+l,t+m+l}_\infty(W)$$
    and
    \begin{enumerate}
        \item $d^{p}_{m}(x)=z$,
        \item $d^{g}_{l}(z)=w$, and has no crossing.
    \end{enumerate}
    Then $d^{q}_{m+l}(x)=w$.
% END AIMPAPER STATEMENT
\uses{paper:2.15}
% Blueprint note: Mathematical dependency review: [AimPaper source, lines 613-615](paper/main.tex.html#L613). Specialize Corollary 2.15 to an identity upper horizontal map.

% Correspondence: fixed original representative-form statement proved.
% Blueprint note: Original statement and context: paper/main.tex.html#L597; paper/2412.10879.pdf.
% Blueprint note: Original-statement interface: \texttt{KIP126.ProofTools.Paper.extensionComposition}. The public statement is proved without changing its type or premises.
% Blueprint note: The printed occurrence of $Y$ for $z$ is read as $Z$, as forced by the diagram; the unused $n$ and the printed $n,m\ge0$ are retained without adding $l\ge0$. See the statement-typing correction record.
% Blueprint note: Apply the lower no-crossing criterion to the actual image of a source representative, then use the composite equality. No additional condition $l\ge0$ is imposed.
\lean{KIP126.ProofTools.Paper.extensionComposition}
\end{corollary}
\begin{proof}
\textit{(See \cite{LWX126}, Corollary 2.17.)} \href{paper/main.tex.html#L613}{Original proof}.

This is a special case of Corollary \ref{cor:0012nik} when $X=Y$ and $f=id_X$.
\end{proof}

\section{Functoriality and exactness of extension spectral sequences}

% AimPaper: 2.18 corollary 617 626
\begin{corollary}[Morphisms of extension spectral sequences]
\label{paper:2.18}
\source{AimPaper Corollary 2.18}
\level{coarse}
From \cite{LWX126}, Corollary 2.18.\par
% BEGIN AIMPAPER STATEMENT
\label{cor:e7b20ae2}
    Consider a homotopy commutative diagram of spectra
    $$\xymatrix{
    X \ar[r]^{f}\ar[d]_{p} & Y\ar[d]^{q}\\
    Z \ar[r]_{g} & W
    }.$$
    Assume $\ESS{p}_0^{*,*}=\ESS{p}_r^{*,*}$ and $\ESS{q}_0^{*,*}=\ESS{q}_r^{*,*}$ for some $r\ge 0$. Then 
    $$(d^p_r, d^q_r): E_\infty^{*,*}(X)\oplus E_\infty^{*,*}(Y)\to E_\infty^{*,*}(Z)\oplus E_\infty^{*,*}(W)$$
    induces a map from the $f$-ESS to the $g$-ESS.
% END AIMPAPER STATEMENT
\uses{paper:2.15,paper:2.1}
% Blueprint note: Mathematical dependency review: [AimPaper source, lines 628-630](paper/main.tex.html#L628). Apply the square corollary inductively to obtain commuting maps on every page of the extension spectral sequences.

% Correspondence: fixed original representative-form statement proved.
% Blueprint note: Original statement and context: paper/main.tex.html#L617; paper/2412.10879.pdf.
% Blueprint note: Original-statement interface: \texttt{KIP126.ProofTools.Paper.extensionSpectralSequenceMorphism}. The public statement is proved without changing its type or premises.
% Blueprint note: Finite induction on the preceding vanishing extension relations proves a filtration shift by $r$. The unique degree-$r$ extension values form additive maps whose exact graphs are the $p$- and $q$-relations; the same homotopy representative and square preserve every $f$-extension relation. This proves the fixed statement in its full representative presentation; the separate construction of a classical ESS remains scoped in Definition 2.1.
\lean{KIP126.ProofTools.Paper.extensionSpectralSequenceMorphism}
\end{corollary}
\begin{proof}
\textit{(See \cite{LWX126}, Corollary 2.18.)} \href{paper/main.tex.html#L628}{Original proof}.

Since $\ESS{p}_0^{*,*}=\ESS{p}_r^{*,*}$ and $\ESS{q}_0^{*,*}=\ESS{q}_r^{*,*}$, we know that $d^p_r$, $d^q_r$ have no crossings. By Corollary \ref{cor:0012nik}, we have $d^g_0\circ d^p_r = d^q_r\circ d^f_0$ and therefore $(d^p_r, d^q_r)$ induces a map from $\ESS{f}_1^{*,*}$ to $\ESS{g}_1^{*,*}$. Inductively we can apply Corollary \ref{cor:0012nik} again and show that $(d^p_r, d^q_r)$ induces a map from $\ESS{f}_n^{*,*}$ to $\ESS{g}_n^{*,*}$ and $d^g_n\circ d^p_r = d^q_r\circ d^f_n$ for all $n$.
\end{proof}

% AimPaper: 2.19 corollary 632 635
\begin{corollary}[Permanent cycles from a null composite]
\label{paper:2.19}
\source{AimPaper Corollary 2.19}
\level{coarse}
From \cite{LWX126}, Corollary 2.19.\par
% BEGIN AIMPAPER STATEMENT
\label{cor:aed3d1a4}
    If the composition $g\circ f$ of two maps $f:X\to Y$, $g:Y\to Z$ is trivial, and $d_n^f(x)=y$ for $x\in E_\infty^{s,t}(X)$ and $y\in E_\infty^{s+n,t+n}(Y)$, then $y$ is a permanent cycle in the $g$-ESS, i.e., we have
    $d^g_m(y)=0$ for all $m\ge 0$.
% END AIMPAPER STATEMENT
\uses{paper:2.15,paper:2.5}
% Blueprint note: Mathematical dependency review: [AimPaper source, lines 637-646](paper/main.tex.html#L637). The null-composite square proves permanence; the article also gives a representative proof using Proposition 2.5.

% Correspondence: fixed original representative-form statement proved.
% Blueprint note: Original statement and context: paper/main.tex.html#L632; paper/2412.10879.pdf.
% Blueprint note: Original-statement interface: \texttt{KIP126.ProofTools.Paper.permanentCycle\_of\_nullComposite}. The public statement is proved without changing its type or premises.
% Blueprint note: The original extension supplies a single representative in the actual image of $f$. Its composite with $g$ is zero, so that same representative witnesses every zero outgoing extension.
\lean{KIP126.ProofTools.Paper.permanentCycle_of_nullComposite}
\end{corollary}
\begin{proof}
\textit{(See \cite{LWX126}, Corollary 2.19.)} \href{paper/main.tex.html#L637}{Original proof}.

By Proposition~\ref{paper:2.5}, choose a representative $[x]$ for which $[y]=f[x]$ is detected by $y$. The null composite gives $g[y]=gf[x]=0$. This same representative therefore has zero image in every filtration quotient. By the representative definition of the extension spectral sequence, every outgoing $g$-extension differential on $y$ vanishes.
\end{proof}

% AimPaper: 2.20 proposition 648 654
\begin{proposition}[Exactness and extension boundaries]
\label{paper:2.20}
\source{AimPaper Proposition 2.20}
\level{coarse}
From \cite{LWX126}, Proposition 2.20.\par
% BEGIN AIMPAPER STATEMENT
\label{prop:cfe810af}
    Suppose that the sequence of spectra
    $$X\fto{f}Y\fto{g}Z$$
    induces a sequence of homotopy groups
    $$\pi_*X\fto{\pi_*f}\pi_*Y\fto{\pi_*g}\pi_*Z$$
    that is exact in $\pi_*Y$. Then all permanent cycles in $E^{*,*}_\infty(Y)$ of the $g$-ESS are boundaries in the $f$-ESS.
% END AIMPAPER STATEMENT
\uses{paper:2.1,paper:2.5,sem:SpectralSequence:completion,sem:AdamsSynthetic:selected-convergence}
% Blueprint note: Mathematical dependency review: [AimPaper source, lines 656-664](paper/main.tex.html#L656). Convergence of the filtered complex and exactness identify permanent cycles with extension boundaries; representatives give the alternative proof.

As throughout \cite{LWX126}, spectra are understood to be 2-completed,
connective and of finite type. The Lean binding expresses the same
conclusion through the representative criterion of
Proposition~\ref{paper:2.5}, with completeness of the middle homotopy
filtration, finite nonnegative filtration quotients and separation of the
target filtration explicit in its background inputs. Permanence still
permits a different representative on each page; zero labels are included.
The fixed Standard model's realization of this completed background and
its convergence hypotheses remains an upstream task, tracked in
issue 31.
% Blueprint note: The AimPaper statement is retained verbatim and directly binds its proved representative-form theorem. The explicit background inputs are retained; this binding does not assert their construction for the fixed Standard model.
% Blueprint note: The deferred completed-background migration is tracked in [issue 31](https://github.com/SII-MATH/Kervaire126Main/issues/31).
\lean{KIP126.ProofTools.Paper.extensionBoundary_of_exact}
\end{proposition}
\begin{proof}
\uses{paper:2.5,sem:SpectralSequence:completion,sem:AdamsSynthetic:selected-convergence}
\textit{(See \cite{LWX126}, Proposition 2.20.)} \href{paper/main.tex.html#L656}{Original proof}.

The representative argument makes the convergence step explicit. Fix
$y\in E_\infty^{s,d+s}(Y)$. For each integer $k$, let $S_k$ consist of the
residue classes in $\pi_dY/F^k\pi_dY$ represented by some $b$ representing
$y$ with $g_*b\in F^k\pi_dZ$. Permanence makes every $S_k$ nonempty.
The finite nonnegative quotients also make every integer-indexed quotient
finite, and the quotient transitions carry $S_l$ into $S_k$ for $k\leq l$.
The finite inverse-system principle gives a compatible family of such
residues. Completeness lifts it to one element $b\in\pi_dY$.

The residue at $k=s+1$ shows that $b$ still represents $y$. For every
$n\geq0$, its residue at $n$ agrees with that of a representative whose
image lies in $F^n\pi_dZ$. Filtration preservation then puts $g_*b$ in
$F^n\pi_dZ$ as well. Separation gives $g_*b=0$, and homotopy exactness
supplies $c\in\pi_dX$ with $f_*c=b$. At $a=\min(0,s)$ the source Adams
filtration is the whole homotopy group, so $c$ defines an
associated-graded class $x$. Proposition~\ref{paper:2.5} gives the required
$f$-extension from $x$ to $y$, so $y$ is an $f$-extension boundary.
\end{proof}

\section{Synthetic boundary maps and extension differentials}

% AimPaper: 4.2 proposition 1025 1029
\begin{proposition}[Extensions for $\lambda^n$ and $\rho$]
\label{paper:4.2}
\source{AimPaper Proposition 4.2}
\level{coarse}
From \cite{LWX126}, Proposition 4.2.\par
% BEGIN AIMPAPER STATEMENT
\label{prop:9770ae6e} 
The only nonzero differentials in the extension spectral sequences for the maps $\lambda^n$ and $\rho$ from Notation~\ref{not:deltaandrho} are the $d_0$'s:
$$d^{\lambda^n}_0=\lambda^n \ \text{and} \ d^{\rho}_0=\rho.$$ 
As a result, these $d_0$'s have no crossings.
% END AIMPAPER STATEMENT
\uses{paper:3.11,paper:3.12,paper:2.19}
% Blueprint note: Mathematical dependency review: [AimPaper source, lines 1030-1050](paper/main.tex.html#L1030). The infinite-page quotient formulas make the lambda/rho sequence exact; the null-composite result rules out higher extensions.

% Correspondence: public counterparts; scope is specified below.
% Blueprint note: The finite-quotient and infinite versions concern the same actual maps $\lambda$ and $\rho$: they exclude essential extensions of positive length, identify the actual $d_0$ equations, and establish the absence of crossings.
% Blueprint note: This entry is bound to the existing Lean declarations. Its status reports their local completion; the scope of correspondence with the article is specified above.
% Blueprint note: Bound declarations (private names are resolved through stable source references without changing their visibility):
% Blueprint note: \begin{itemize}
% Blueprint note: \item \texttt{KIP126.ProofTools.Paper.lambda\_rho\_extension\_collapse\_finite}
% Blueprint note: \item \texttt{KIP126.ProofTools.Paper.lambda\_rho\_extension\_collapse\_infinite}
% Blueprint note: \end{itemize}
\lean{KIP126.ProofTools.Paper.lambda_rho_extension_collapse_finite,KIP126.ProofTools.Paper.lambda_rho_extension_collapse_infinite}

% Blueprint note: The original article number is 4.2; other numbering is generated by the blueprint. Throughout the article, spectra are assumed to be 2-completed, connective, and of finite type. Read this statement together with its original context.
% Blueprint note: \href{paper/main.tex.html#L1025}{Original statement, hypotheses, and context} · \href{paper/2412.10879.pdf}{Article PDF}
\end{proposition}
\begin{proof}
\textit{(See \cite{LWX126}, Proposition 4.2.)} \href{paper/main.tex.html#L1030}{Original proof}.

Compute the three associated graded groups in the cofiber sequence using Propositions~\ref{paper:3.11} and~\ref{paper:3.12}. Write $k=t-w$. Outside $0\le k<m$ the middle group vanishes. For $0\le k<n$, the $\lambda^n$ source group vanishes and $\rho$ is the inclusion of the relevant cycle subquotients. For $n\le k<m$, the $\rho$ target group vanishes and $\lambda^n$ is the quotient map from a smaller boundary quotient onto the middle group. Thus the sequence of associated graded groups is exact in the middle in every degree.

The two composites are zero. Corollary~\ref{paper:2.19} makes the image of each preceding map a permanent cycle for the next extension spectral sequence. Since each nonzero graded $\lambda^n$ map is surjective and each nonzero graded $\rho$ map is injective in these ranges, no positive-length extension remains. Hence the only differentials are the indicated $d_0$ maps, which have no crossings. The untruncated case uses $Z_\infty$ in the same calculation.
\end{proof}

For $n>0$, the absolute connecting map of the quotient triangle is
\[\delta_n:\nu X/\lambda^n\longrightarrow\Sigma^{1,-n}\nu X.\]
For $0<n<m$, its finite version has target $\Sigma^{1,-n}(\nu X/\lambda^{m-n})$. Using the suspension isomorphism of the quotient triangle, the source label of $x$ is $\lambda^a x$ in the $n$-th quotient. For a classical differential of length $r=k+1$ and $n\leq a+k$, its target label is the shifted $\lambda^{a+k-n}y$; the finite label is its image under the target quotient projection. A crossing of the absolute boundary on page $p$ is an essential extension in that same boundary ESS, with source filtration $u$ and target filtration $v$ satisfying $s<u\leq v\leq s+r$.

% AimPaper: 4.6 proposition 1068 1086
\begin{proposition}[Boundary extensions from Adams differentials]
\label{paper:4.6}
\source{AimPaper Proposition 4.6}
\level{coarse}
From \cite{LWX126}, Proposition 4.6.\par
% BEGIN AIMPAPER STATEMENT
\label{prop:6de7d130}
    Suppose in the classical Adams spectral sequence of $X$ we have $d_r(x)=y$, where $x\in Z_{r-1}^{s,t}(X)$ and $y\in Z_\infty^{s+r,t+r-1}(X)$. Consider the map
    $$\delta:\nu X/\lambda^n\to \Sus^{1,-n}\nu X/\lambda^{m-n}.$$
    \begin{enumerate}
        \item If $r\ge n+1$, then we view $x$ as an element of 
        $$E_\infty^{s,t,t}(\nu X/\lambda^n)\iso Z_n^{s,t}(X),$$
        and $\lambda^{r-n-1} y$ as an element of 
        $$E_\infty^{s+r,t+r-1, t+n}(\nu X/\lambda^{m-n})\iso Z_{m-r+1}^{s+r,t+r-1}(X)/B_{r-n}^{s+r,t+r-1}(X).$$
        We then have
        $$d^{\delta}_r(x)=\lambda^{r-n-1}y,$$
        which is trivial if $r>m$.
        \item If $r<n+1$, then we view $\lambda^{n+1-r}x$ as an element of 
        $$E_\infty^{s,t,t-n-1+r}(\nu X/\lambda^n)\iso Z^{s,t}_{r-1}(X)/B^{s,t}_{n+2-r}(X),$$
        and $y$ as an element of 
        $$E_\infty^{s+r,t+r-1, t+r-1}(\nu X)\iso Z_\infty^{s+r,t+r-1}(X).$$
        In this case, we have
        $$d^{\delta}_r(\lambda^{n+1-r}x)=y.$$
    \end{enumerate}
% END AIMPAPER STATEMENT
\uses{paper:3.8,paper:3.11,paper:3.12,paper:4.2,paper:2.16,paper:2.17}
% Blueprint note: Mathematical dependency review: [AimPaper source, lines 1087-1131](paper/main.tex.html#L1087). Start with the Bockstein comparison, induct on the quotient length using lambda and restriction triangles, and reduce the finite target by composition of extensions.

% Correspondence: public counterparts; scope is specified below.
% Blueprint note: The four public statements cover both ranges of differential lengths, for finite and absolute $\delta$-extensions.
% Blueprint note: This entry is bound to the existing Lean declarations. Its status reports their local completion; the scope of correspondence with the article is specified above.
% Blueprint note: Bound declarations (private names are resolved through stable source references without changing their visibility):
% Blueprint note: \begin{itemize}
% Blueprint note: \item \texttt{KIP126.ProofTools.Paper.syntheticDeltaExtension\_prop46\_above\_absolute}
% Blueprint note: \item \texttt{KIP126.ProofTools.Paper.syntheticDeltaExtension\_prop46\_below\_absolute}
% Blueprint note: \item \texttt{KIP126.ProofTools.Paper.syntheticDeltaExtension\_prop46\_above\_finite}
% Blueprint note: \item \texttt{KIP126.ProofTools.Paper.syntheticDeltaExtension\_prop46\_below\_finite}
% Blueprint note: \end{itemize}
\lean{KIP126.ProofTools.Paper.syntheticDeltaExtension_prop46_above_absolute,KIP126.ProofTools.Paper.syntheticDeltaExtension_prop46_below_absolute,KIP126.ProofTools.Paper.syntheticDeltaExtension_prop46_above_finite,KIP126.ProofTools.Paper.syntheticDeltaExtension_prop46_below_finite}

% Blueprint note: The original article number is 4.6; other numbering is generated by the blueprint. Throughout the article, spectra are assumed to be 2-completed, connective, and of finite type. Read this statement together with its original context.
% Blueprint note: \href{paper/main.tex.html#L1068}{Original statement, hypotheses, and context} · \href{paper/2412.10879.pdf}{Article PDF}
\end{proposition}
\begin{proof}
\textit{(See \cite{LWX126}, Proposition 4.6.)} \href{paper/main.tex.html#L1087}{Original proof}.

First take $m=\infty$ and write $\delta_n=\delta_{n,\infty}$. For $n=1$, Theorem~\ref{paper:3.8} identifies the connecting-map extension with the classical Adams differential. Induct on $n$, using the triangle comparing the quotients of lengths $n-1$ and $n$ and multiplication by $\lambda$.

For $r\le n$, the commutative-triangle extension rule transports the induction hypothesis to $\lambda^{n+1-r}x$ and gives target $y$. For $r\ge n+1$, the comparison gives $\lambda d_r^{\delta_n}(x)=\lambda^{r-n}y$ modulo classical $B_{r-1}$. The possible ambiguity in dividing this particular equality by $\lambda$ is contained in $B_{r-n+1}\subseteq B_{r-1}$, by the explicit infinity-page subquotients. This yields $d_r^{\delta_n}(x)=\lambda^{r-n-1}y$ in the required quotient, rather than an unrestricted cancellation theorem.

Induction on the differential length identifies $B_{r-1}$ here with the sum of the earlier $\delta_n$-extension images: each earlier classical boundary gives its corresponding boundary extension, and the same comparison gives the converse. Finally $\delta_{n,m}$ is $\delta_n$ followed by restriction to the finite target quotient. Proposition~\ref{paper:4.2} and the composition rule transport both formulas to that quotient. The target is zero when its $\lambda$-exponent reaches the quotient length, giving the stated vanishing range.
\end{proof}

% AimPaper: 4.10 corollary 1138 1144
\begin{corollary}[$\lambda$ multiples of boundary extensions]
\label{paper:4.10}
\source{AimPaper Corollary 4.10}
\level{coarse}
From \cite{LWX126}, Corollary 4.10.\par
% BEGIN AIMPAPER STATEMENT
\label{cor:2a636737}
    For $x,y,\delta$ in Proposition \ref{prop:6de7d130} we always have
    \begin{equation}\label{eq:841425e9}
        d^{\delta}_r(\lambda^ax)=\lambda^{a+r-n-1}y
    \end{equation}
    if $0\le a\le n$ and $0\le a+r-n-1<m-n$ (the differential is trivial if $a$ exceeds this range).
% END AIMPAPER STATEMENT
\uses{paper:4.6,paper:4.2}
% Blueprint note: Mathematical dependency review: [AimPaper source, lines 1138-1150](paper/main.tex.html#L1138). Multiply the boundary-extension formulas by lambda; the following remark identifies the complete classical differential coset.

% Correspondence: public counterparts; scope is specified below.
% Blueprint note: The public statements cover the full valid range of weights, the zero endpoint $a=n$, and vanishing when the source or target lies outside the strip.
% Blueprint note: This entry is bound to the existing Lean declarations. Its status reports their local completion; the scope of correspondence with the article is specified above.
% Blueprint note: Bound declarations (private names are resolved through stable source references without changing their visibility):
% Blueprint note: \begin{itemize}
% Blueprint note: \item \texttt{KIP126.ProofTools.Paper.syntheticDeltaExtension\_absolute}
% Blueprint note: \item \texttt{KIP126.ProofTools.Paper.syntheticDeltaExtension\_finite}
% Blueprint note: \item \texttt{KIP126.ProofTools.Paper.syntheticDeltaSource\_vanishes}
% Blueprint note: \item \texttt{KIP126.ProofTools.Paper.syntheticDeltaFiniteTarget\_vanishes}
% Blueprint note: \end{itemize}
\lean{KIP126.ProofTools.Paper.syntheticDeltaExtension_absolute,KIP126.ProofTools.Paper.syntheticDeltaExtension_finite,KIP126.ProofTools.Paper.syntheticDeltaSource_vanishes,KIP126.ProofTools.Paper.syntheticDeltaFiniteTarget_vanishes}

% Blueprint note: The original article number is 4.10; other numbering is generated by the blueprint. Throughout the article, spectra are assumed to be 2-completed, connective, and of finite type. Read this statement together with its original context.
% Blueprint note: \href{paper/main.tex.html#L1138}{Original statement, hypotheses, and context} · \href{paper/2412.10879.pdf}{Article PDF}
\end{corollary}
\begin{proof}
\textit{(See \cite{LWX126}, Corollary 4.10.)} \href{paper/main.tex.html#L1138}{Original statement and context}.

Apply multiplication by $\lambda^a$ to the boundary-extension formulas of Proposition~\ref{paper:4.6}. Naturality of the quotient cofiber sequence makes this multiplication commute with the connecting map; Proposition~\ref{paper:4.2} ensures that it introduces no additional positive-length extension. The source label is multiplied by $\lambda^a$ and the target label becomes $\lambda^{a+r-n-1}y$.

The two labels must lie in the source and target windows of Proposition~\ref{paper:3.12}. These are exactly the ranges displayed in the statement. Beyond those windows the relevant layer vanishes, so the corresponding extension is trivial. The quotient interpretation of the labels, including earlier boundaries, is the one fixed in Proposition~\ref{paper:4.6}.
\end{proof}

\section{Crossings of Adams differentials}

% AimPaper: 4.14 definition 1165 1173
\begin{definition}[Crossings of Adams differentials]
\label{paper:4.14}
\source{AimPaper Definition 4.14}
\level{coarse}
From \cite{LWX126}, Definition 4.14.\par
% BEGIN AIMPAPER STATEMENT
\label{def:classicalcrossdiff}
    Suppose $r\ge n+1$ and $d_r(x)=y$, where
    $$x\in E_2^{s,t}(X), \hspace{0.5cm} y\in E_2^{s+r,t+r-1}(X).$$
    A crossing of $d_r(x)=y$ on the $E_{n+1}$-page refers to an essential Adams differential
    $$d_{r-a-b}(x^\prime)=y^\prime,$$
    where
    $$x^\prime\in E_2^{s+a,t+a}(X), \hspace{.5cm} y^\prime\in E_2^{s+r-b,t+r-b-1}(X),$$
    with $0<a\le n-1$ and $0\le b\le r-n-1$. See Figure \ref{fig:f5eae42a}.
% END AIMPAPER STATEMENT
\uses{sem:SpectralSequence:differential,sem:SpectralSequence:cycles-boundaries}
% Blueprint note: Mathematical dependency review: [AimPaper source, lines 1165-1173](paper/main.tex.html#L1165). Define a classical crossing by a shorter essential Adams differential in the specified filtration rectangle.

% Correspondence: public counterparts; scope is specified below.
% Blueprint note: The bound definition uses the direction of classical differential crossings and the valid page range specified in the article. It is distinct from Moss crossings. The separate no-crossing predicate is not defined in this entry and is not bound here.
\lean{KIP126.ProofTools.HasCrossingOn}

% Blueprint note: The original article number is 4.14; other numbering is generated by the blueprint. Throughout the article, spectra are assumed to be 2-completed, connective, and of finite type. Read this statement together with its original context.
% Blueprint note: \href{paper/main.tex.html#L1165}{Original statement, hypotheses, and context} · \href{paper/2412.10879.pdf}{Article PDF}
\end{definition}

% AimPaper: 4.16 proposition 1218 1224
\begin{proposition}[Adams crossings and boundary-extension crossings]
\label{paper:4.16}
\source{AimPaper Proposition 4.16}
\level{coarse}
From \cite{LWX126}, Proposition 4.16.\par
% BEGIN AIMPAPER STATEMENT
\label{prop:cross-dr-En}
    The Adams differential $d_r(x)=y$ has a crossing on the $E_{n+1}$-page if and only if the corresponding $\delta_n$-extension
    $$d^{\delta_n}_r(x)=\lambda^{r-n-1}y$$
    for
    $$\delta_{n}:\nu X/\lambda^{n}\to \Sus^{1,-n}\nu X$$
    has a crossing.
% END AIMPAPER STATEMENT
\uses{paper:4.14,paper:2.7,paper:3.11,paper:3.12,paper:4.6,paper:4.10}
% Blueprint note: Mathematical dependency review: [AimPaper source, lines 1225-1234](paper/main.tex.html#L1225). Translate extension crossings through the infinity-page quotient coordinates and the lambda-multiple boundary formula.

% Correspondence: public counterparts; scope is specified below.
% Blueprint note: The public theorem gives the equivalence between classical crossings and crossings in the actual $\delta$-extension spectral sequence throughout the range $2\leq\mathrm{page}\leq r$.
% Blueprint note: This entry is bound to the existing Lean declarations. Its status reports their local completion; the scope of correspondence with the article is specified above.
% Blueprint note: Bound declarations (private names are resolved through stable source references without changing their visibility):
% Blueprint note: \begin{itemize}
% Blueprint note: \item \texttt{KIP126.ProofTools.Paper.classical\_crossing\_iff\_delta\_crossing}
% Blueprint note: \end{itemize}
\lean{KIP126.ProofTools.Paper.classical_crossing_iff_delta_crossing}

% Blueprint note: The original article number is 4.16; other numbering is generated by the blueprint. Throughout the article, spectra are assumed to be 2-completed, connective, and of finite type. Read this statement together with its original context.
% Blueprint note: \href{paper/main.tex.html#L1218}{Original statement, hypotheses, and context} · \href{paper/2412.10879.pdf}{Article PDF}
\end{proposition}
\begin{proof}
\textit{(See \cite{LWX126}, Proposition 4.16.)} \href{paper/main.tex.html#L1225}{Original proof}.

By Propositions \ref{prop:30e8b746}, \ref{prop:59f111f} and \ref{prop:6de7d130}, a crossing of $d^{\delta_n}_r(x)=\lambda^{r-n-1}y$ takes the form
    $$d^{\delta_n}_{r-a-b}(\lambda^a x')=\lambda^{r-n-1-b}y',$$
    where $0<a\le n-1$, $0\le b\le r-n-1$,
    $$\lambda^a x'\in E_\infty^{s+a,t+a,t}(\nu X/\lambda^n)\iso Z_{n-a}^{s+a,t+a}(X)/B_{1+a}^{s+a,t+a}(X),$$
    and
    $$y'\in E_\infty^{s+r-b, t+r-b-1}(\nu X).$$
    By Corollary \ref{cor:2a636737}, we see that this crossing corresponds to the classical Adams differential
    $$d_{r-a-b}(x')=y'.$$
\end{proof}

\section{Page extensions at zero and positive lambda levels}

\begin{definition}[Level extensions and the quotient-boundary square]
\label{sem:ProofTools:level-finite-extension}
\level{coarse}
Following the arbitrary-level quotient windows preceding Definition 5.4 of \cite{LWX126}, fix a quotient exponent $q$, a normalization shift $e$, a level $k\in\mathbb N$, and $x\in E_2^{s,t}(X)$. The normalized finite source label is the image of the native label $\lambda^k x$ under the actual $(0,e)$ shift and the inverse quotient-shift comparison, in tridegree
\[(s,t,t+e-k)\quad\text{of}\quad (\Sigma^{0,e}\nu X)/\lambda^q.\]
The chosen comparison respects both quotient inclusion and absolute boundary. This specifies an actual source label without assuming a new calibration of the existing comparison at positive shift and positive level.

Let $f:X\to Y$, $e=e(f)$, $A=\Sigma^{0,e}\nu X$, $B=\nu Y$, and $F:A\to B$ be the normalized map. For $n,k\in\mathbb N$, put $q=n+k-1$. A level-$k$ finite extension of length $m$ from $x\in E_2^{s,t}(X)$ to $y\in E_2^{s+m,t+m}(Y)$ consists of the bounds
\[2\le n,\qquad e\le m\le n-2+e,\]
and actual infinite-page source and target classes, represented by the normalized label $\lambda^k x$ and the quotient projection of $\lambda^{k+m-e}y$, which satisfy
\[d_m^{F_q}(\lambda^k x)=\lambda^{k+m-e}y\]
at fixed stem and weight $(t-s,t+e-k)$. The respective native windows are
\[Z_{n-1}^{s,t}(X)/B_{k+1}^{s,t}(X),\qquad
Z_{n-1-m+e}^{s+m,t+m}(Y)/B_{k+m-e+1}^{s+m,t+m}(Y).\]
The extension need not be essential, and zero labels are allowed.
We abbreviate it by $\mathcal E^f_{q,k}(m;x,y)$, with the coordinate
convention of Definition~\ref{sem:AdamsSynthetic:page-extension-levels}.
Its complete target indeterminacy includes the shorter images of this
same finite map; it is not determined by nonvanishing of the chosen
classical target alone.

For an actual filtered map $g$ and fixed stem/weight $p$, write $NC(g;p,b,h)$ if there is no essential extension $d_{v-u}^g(\xi)=\eta$ for any $b<u\le v\le b+h$ and any associated-graded classes $\xi,\eta$. Essentiality is taken modulo the complete target indeterminacy of that actual map. For the further lengths $r,l$, with $A,B,F,q$ as above and $p=(t-s,t+e-k)$, the three predicates are respectively
\[
NC(F_q;p,s,m),\qquad NC(\delta_q^A;p,s,r),\qquad
NC(\Sigma^{1,-q}F;p,s+r,l).
\]
They refer to the top, left and bottom maps in the quotient-boundary square. The bottom condition is tested on the propagated infinite extension, at its actual level $r-n$; it is not inferred from a zero-level no-crossing condition.
\uses{sem:AdamsSynthetic:page-extension-levels,sem:AdamsSynthetic:ess-shifts,sem:ProofTools:mh-actual-no-crossing}
\lean{KIP126.ProofTools.Standard.levelFiniteSourceLabel,KIP126.ProofTools.Standard.LevelFiniteExtension,KIP126.ProofTools.Standard.LevelFiniteNoCrossing,KIP126.ProofTools.Standard.LevelBoundaryNoCrossing,KIP126.ProofTools.Standard.LevelBottomNoCrossing}
\end{definition}


\begin{theorem}[From a page extension to its level-zero representative form]
\label{sem:ProofTools:extension-level-zero}
\level{coarse}
A representative comparison for Definition 5.4 of \cite{LWX126} in the
fixed model.\par
Let $f:X\to Y$, $n\in\mathbb N$ with $n\geq2$, and
$m,s,t\in\mathbb Z$. Suppose that the specified labels
$x\in E_2^{s,t}(X)$ and $y\in E_2^{s+m,t+m}(Y)$ support the
finite extension $d_m^{f,E_n}(x)=y$. Then they support
\[
\mathcal E^f_{n-1,0}(m;x,y).
\]
The same actual map equation and source and target labels are retained.
No additional cycle or crossing premise is required. This assertion
does not identify the two sets of no-crossing hypotheses.
\uses{paper:5.4,sem:ProofTools:level-finite-extension}
\lean{KIP126.ProofTools.Standard.finiteExtensionAtLevelZero}
% Correspondence: the bound result proves the stated zero-level implication. A converse, arbitrary-level stabilization and transport of crossing predicates are not asserted by this entry.
\end{theorem}
\begin{proof}
Choose the actual source and target witnesses in the page extension.
The zero-level shift and quotient comparisons identify their labels
with the level-zero labels in the preceding definition. Transporting
these two representative equalities retains the same extension equation.
The page and length bounds are already part of the given witness.
\end{proof}

Fixing $n$ and increasing $k$ replaces the quotient exponent $n-1$
by $n+k-1$ and the weight by $t+e-k$. This is different from quotient
reduction at fixed weight, and different again from division by
$\lambda$. Each operation must retain its specified labels before it
can be used in an extension argument. Neither this zero-level comparison
nor naturality alone supplies preservation of essentiality or actual
no-crossing at arbitrary levels.

\subsection{Examples of page extensions}

The following source examples illustrate the page-extension language.
They are not hypotheses of the general tools. Their derivations appear
with the corresponding Mahowald and stretching applications below.
Existence of a relation, essentiality, and absence of crossing are
separate assertions for each specified map and page.

\begin{proposition}[Page extensions for nu and multiplication by two]
\label{paper:5.8}
\label{exam:extonEn}
\level{coarse}
From \cite{LWX126}, Example 5.8.\par
For $f=\nu:S^3\to S^0$, with $e(f)=1$, the article gives the following extensions:
\[
\begin{aligned}
d_1^{f,E_2}(h_5)&=h_2h_5,\\
d_2^{f,E_3}(h_0h_4^2)&=h_0p,\\
d_2^{f,E_\infty}(h_0h_4^2)&=h_0p.
\end{aligned}
\]
The $E_3$ relation means that a representative of $h_0h_4^2$ in $S^{3,4}/\lambda^2$ has $[h_2]$-image detecting $\lambda h_0p$ in $S^{0,0}/\lambda^2$. The $E_\infty$ relation requires representatives in the full synthetic spheres with the same specified source and image detections.

For $f=2:S^0\to S^0$, normalized by $\lambda[h_0]=2$, the article gives
\[
d_2^{f,E_\infty}(h_0h_3^2)=0,
\qquad d_1^{f,E_\infty}(d_0)=h_0d_0.
\]
The latter extension is essential and makes
\[
d_2^{f,E_\infty}(h_0h_3^2)=h_0d_0
\]
inessential. Thus two target labels can describe the same full extension coset.
\uses{paper:5.4,paper:5.7,sem:Literature:hopf-classes,sem:Literature:synthetic-h-zero}
\lean{KIP126.ProofTools.Standard.paper5_8_nu_e2_extension,KIP126.ProofTools.Standard.paper6_19_finite_nu_extension,KIP126.ProofTools.Standard.paper6_24_infinite_nu_extension,KIP126.ProofTools.Standard.paper6_8_complete_crossing_warning}
% Correspondence: The four public results cover the specified E2 nu extension, its E3 and infinite numerical extensions, and the complete multiplication-by-two assertions. The E2 source is the actual threefold suspended h5 label; the full nu and multiplication-by-two representatives are constructed independently of finite-quotient restriction. The positive-level questions below are not additional conclusions of these bindings.
\end{proposition}
\begin{proof}
\uses{paper:6.19,paper:6.24}
The first relation is the $E_2$ product with the detector $h_2$ of $\nu$, using the actual threefold suspension of $h_5$ and the normalized map on the first quotient. The second relation is the finite nu extension proved in Example~\ref{paper:6.19}. For the third, use the actual full representatives constructed in Example~\ref{paper:6.24}, with their specified source and target detections and strict normalized-map equation. The infinite relation is thereby established without inferring a full lift from finite-quotient equations alone.

For multiplication by two, the full representative calculation in Example~\ref{paper:6.8} gives a representative of $h_0h_3^2$ annihilated by $[h_0]$. The shorter essential extension from $d_0$ to $h_0d_0$ puts the latter label in the length-two target indeterminacy. Adding its representative constructs the displayed nonzero-target length-two relation and proves its inessentiality.

\textit{The level-indexed questions for the four examples.}
For $a\geq0$, the finite versions of (1) and (2) ask respectively for
\[
\mathcal E^\nu_{a+1,a}(1;h_5,h_2h_5),\qquad
\mathcal E^\nu_{a+2,a}(2;h_0h_4^2,h_0p).
\]
The suspension of the source labels is the same as in the original
$\nu:S^3\to S^0$ examples. The infinite versions in (3) and (4)
instead ask for representatives in the full normalized map, with
source label multiplied by $\lambda^a$ and target label multiplied
by $\lambda^{a+n-e(f)}$. Such equations require their actual
representatives; equations in each finite quotient separately do not
supply them. In (4), whether the transported shorter extension
remains essential is an additional target-coset calculation.
These are the precise level questions for the numerical claims,
not an assertion that all of their inputs have been formally produced.
\end{proof}

\section{Crossing comparisons and conditional lambda division}

The classical candidate definition and the actual no-crossing predicates
have different quantifiers. The conditional division theorem below
replaces the division step in the original crossing comparison while
preserving actual source and image data. The original comparison and
its counterexamples are retained in the article and the AimPaper corrections.

% AimPaper: 5.9 definition 1437 1441

\begin{definition}[Crossings of $E_r$-extensions]
\label{paper:5.9}
\source{AimPaper Definition 5.9}
\level{coarse}
From \cite{LWX126}, Definition 5.9.\par
% BEGIN AIMPAPER STATEMENT
\label{def:41d51149}
    A crossing of the $(f,E_r)$-extension $d_{n}^{f,E_r}(x)=y$ in (\ref{eq:b98be76a}) is defined as an essential $(f, E_{r-a})$-extension from some $x'\in Z_{r-1-a}^{s+a,t+a}(X)$ to 
    $$y'\in Z_{r-1-n+b+e(f)}^{s+n-b,t+n-b}(Y)\backslash B_{1+n-b-e(f)}^{s+n-b,t+n-b}(Y)$$
    (where $\backslash$ denotes the difference of sets and it means that $y'$ should survive to the classical Adams $E_{r-n+b+e(f)}$ page while it should not be hit by an Adams differential of length at most $1+n-b-e(f)$) for $0<a \le r-2$ and $0\le b\le n-a-e(f)$. See Figure \ref{fig:66241bf6}.
% END AIMPAPER STATEMENT
\uses{paper:5.4}
% Blueprint note: Mathematical dependency review: [AimPaper source, lines 1437-1441](paper/main.tex.html#L1437). Define a crossing by an essential page extension with the displayed smaller-page and filtration bounds.

% Correspondence: public counterparts; scope is specified below.
% Blueprint note: The binding describes the finite-page crossing defined above. The separate infinite-page and no-crossing predicates are not defined in this entry and are not bound here.
\lean{KIP126.ProofTools.FiniteHasCrossing}

% Blueprint note: The original article number is 5.9; other numbering is generated by the blueprint. Throughout the article, spectra are assumed to be 2-completed, connective, and of finite type. Read this statement together with its original context.
% Blueprint note: \href{paper/main.tex.html#L1437}{Original statement, hypotheses, and context} · \href{paper/2412.10879.pdf}{Article PDF}
\end{definition}

% AimPaper: 5.10 proposition 1515 1517

\begin{theorem}[Conditional finite lambda division preserving the actual source and image]
\label{sem:ProofTools:finite-lambda-division}
\source{A conditional replacement for the division step in AimPaper Proposition 5.10}
\level{coarse}
A conditional finite-quotient replacement for the lambda-division step in
\cite{LWX126}, Proposition 5.10.\par
The original proposition asserts an equivalence between classical and
synthetic crossings. The theorem below supplies the required division
operation under an additional target-kernel crossing condition; it does
not assert that equivalence.

Let $f:X\to Y$, $e=e(f)$, $A=\Sigma^{0,e}\nu X$, $B=\nu Y$, and $F:A\to B$ be the normalized map. Fix $0<i<q$, $p=(t-s,t+e)$ and a target filtration $N$. In the actual lambda square put
\[
A'=\Sigma^{0,-i}(A/\lambda^{q-i}),\quad
B'=\Sigma^{0,-i}(B/\lambda^{q-i}),\quad
F'=\Sigma^{0,-i}F_{q-i}.
\]
Let $L_A:A'\to A/\lambda^q$ and $L_B:B'\to B/\lambda^q$ be the actual finite lambda maps. The extra crossing condition states that, for every
\[s+i\le u\le v<N,\]
an essential extension $d_{v-u}^{F'}(\xi)=\eta$ at $p$ has target satisfying
\[\operatorname{gr}_v(L_B)(\eta)\ne0.\]
The initial source row $u=s+i$ is included, while the terminal target row $v=N$ is excluded. This rules out essential smaller-layer targets killed by lambda; it neither requires injectivity on the entire graded target nor follows merely from absence of positive-length extensions of the lambda map itself.

Assume $s+i\le N$ and the additional target-kernel crossing condition. If an actual representative
\[a\in F^{s+i}\pi_{t-s,t+e}(A/\lambda^q)\]
has $F_q(a)\in F^N\pi_{t-s,t+e}(B/\lambda^q)$, there exists
\[b\in F^{s+i}\pi_{t-s,t+e}(A')\]
such that
\[L_A(b)=a,\qquad L_B(F'(b))=F_q(a),\qquad F'(b)\in F^N\pi_{t-s,t+e}(B').\]
Thus division preserves the specified actual source and image and reflects the requested target filtration. The smaller-layer source class is the class of the constructed $b$; no independently selected smaller source is prescribed, and kernel ambiguity at the final target grade is allowed.
\uses{sem:ProofTools:finite-correction,sem:ProofTools:finite-strip,sem:AdamsSynthetic:lambda-quotients}
\lean{KIP126.ProofTools.Standard.finiteLambdaDivision}
% Blueprint note: The extra condition includes the initial source row and excludes the terminal target row. This sufficient-condition theorem does not prove the original Proposition 5.10 equivalence. The original public comparison and its registered proof hole remain in Lean, without a standalone Blueprint node. Registered upstream obligations remain.
\end{theorem}
\begin{proof}
\textit{Proof sketch in the fixed model.}
Naturality of the lambda square alone does not preserve the required
image filtration after division: low-filtration $\lambda$-torsion can
occur. The extra target-kernel crossing condition excludes precisely
the essential lower-filtration images that obstruct the following argument.
Lift the actual source along the finite $\lambda^i$ map. Correct this lift over the finite source and target filtration intervals. A first lower target grade killed by the large-quotient lambda map would give an essential smaller-layer extension forbidden by the stated kernel condition. Finite correction therefore raises the smaller-layer image to $F^N$ while retaining the actual large-layer source and image. The two displayed map equations follow from the actual lambda square. No predetermined smaller source label or injectivity at the terminal target grade is needed.
\end{proof}

% Native hgraph authoring; revised entries retain their original AimPaper source locators.
% Source locators for the corrections tab; no Blueprint statements are retained.
% AimPaper: 6.1 theorem 1606 1621
% AimPaper: 6.12 theorem 1780 1799

\subsection{Crossings in the numerical examples}

\begin{proposition}[Crossings in the nu and multiplication-by-two examples]
\label{paper:5.13}
\label{exam:Ercross}
\level{coarse}
From \cite{LWX126}, Example 5.13.\par
The following are the crossing assertions for the maps and labels of Example~\ref{paper:5.8}, using the classical definition in \ref{paper:5.9}. Assume the displayed extension relations exist; in the fourth case assume the displayed length-one extension is essential. The assertions concern their crossings, rather than establishing existence of the input extensions:
\begin{enumerate}
\item For $f=\nu:S^3\to S^0$, $d_1^{f,E_2}(h_5)=h_2h_5$ has no crossing.
\item For the same $f$, $d_2^{f,E_3}(h_0h_4^2)=h_0p$ has no crossing.
\item The corresponding $d_2^{f,E_\infty}(h_0h_4^2)=h_0p$ has no crossing.
\item For $f=2:S^0\to S^0$, the essential extension $d_1^{f,E_\infty}(d_0)=h_0d_0$ crosses both the length-two relation with target zero and the inessential length-two relation with target $h_0d_0$.
\end{enumerate}
These are assertions about the displayed maps, labels and pages, not claims that no-crossing automatically persists at every lambda level.
\uses{paper:5.9}
\lean{KIP126.ProofTools.Standard.paper5_13_nu_e2_no_crossing,KIP126.ProofTools.Standard.paper5_13_nu_e3_no_crossing,KIP126.ProofTools.Standard.paper5_13_nu_infinite_no_crossing,KIP126.ProofTools.Standard.paper6_8_complete_crossing_warning}
% Correspondence: The three public nu theorems exclude the actual finite and infinite crossing witnesses in the stated positions; the multiplication-by-two theorem supplies the essential short crossing and both long target presentations. These conclusions do not assert actual synthetic no-crossing in another quotient, weight or level and do not use the original Proposition 5.10.
\end{proposition}
\begin{proof}
The first relation is in the minimal-length range: its crossing definition would require $0<a\leq0$. For the second and third relations, the length bounds leave only $a=1,b=0$. The corresponding source group is generated by $h_0^2h_4^2$. The actual first-quotient normalized nu operation is multiplication by $h_2$, and $h_2h_0^2h_4^2=0$; hence that operation is zero on the entire candidate source group. In the finite case every possible image is zero, so an essential candidate is impossible. In the full case a zero first-quotient image lies one step higher in the target filtration, again contradicting essentiality. The fourth assertion follows from the constructed shorter essential extension in Example~\ref{paper:6.8} and its length-two target indeterminacy.

The alternative synthetic argument printed in the article invokes
\cite{LWX126}, Proposition 5.10. Its unrestricted crossing comparison
has the division issue addressed by the extra hypothesis of
Theorem~\ref{sem:ProofTools:finite-lambda-division}. Consequently this
example does not supply an unproved conversion of the classical checks
into the actual no-crossing hypotheses of the level-indexed Leibniz theorem.

At level $a$, the first two finite examples have quotient exponents
$a+1$ and $a+2$. Their actual crossing questions concern the maps
and weights of those quotients. The two infinite examples concern
the full map at the shifted weight. A transported short relation
need not remain essential, and disappearance of one target under
$\lambda$ does not exclude every other essential crossing.
The original four checks therefore neither prove all-level
no-crossing nor constitute a new general crossing comparison.
\end{proof}

\section{Classical realization and target-preserving lifting}

A full synthetic extension can be transported to a classical homotopy
extension using the following realization comparison. Producing a full
representative from finite data is a separate lifting problem, described
after that comparison. These lifting requirements belong to classical
extension applications; they are not extra premises of the finite
composition theorem developed later.

% AimPaper: 5.14 proposition 1581 1587

\begin{proposition}[From infinity-page extensions to homotopy extensions]
\label{paper:5.14}
\source{AimPaper Proposition 5.14}
\level{coarse}
From \cite{LWX126}, Proposition 5.14.\par
% BEGIN AIMPAPER STATEMENT
Consider $f: X\to Y$, $x\in Z_\infty^{s,t}(X)$ and $y\in Z_\infty^{s+n,t+n}(Y)$.
    Then
    $$d^{f,E_\infty}_n(x)=y$$
    implies that
    $$d^{f}_n(x+B_\infty^{s,t}(X))=y+B_\infty^{s,t}(Y).$$
% END AIMPAPER STATEMENT
\uses{paper:5.4,sem:Literature:lambda-inversion,paper:3.11,paper:2.5,sem:Literature:realization-coordinates,sem:Literature:lambda-inversion-adams-comparison,sem:Literature:lambda-inversion-detection}
% Blueprint note: Mathematical dependency review: [AimPaper source, lines 1589-1597](paper/main.tex.html#L1589). Inverting lambda gives a map from the synthetic ESS to the classical ESS and sends the infinity-page extension to its homotopy extension.

As throughout \cite{LWX126}, spectra are understood to be 2-completed,
connective and of finite type. The Lean binding expresses the extension
equation by the representative criterion of Proposition~\ref{paper:2.5}:
$x$ detects one $a$ and $y$ detects $f_*a$. Realization coordinates and
the actual Adams-tower comparison of
Definition~\ref{sem:Literature:lambda-inversion-adams-comparison} make the
lambda-inversion background explicit. Zero labels, boundary labels and
inessential extensions remain allowed. Constructing these comparison
data for the fixed Standard model remains an upstream task, tracked in
issue 31.

\textit{Positive levels and finite quotients.}
A full-object equation at level $k$ uses source $\lambda^kx$
and target $\lambda^{k+n-e}y$ at weight $t+e-k$.
Comparing it with the classical equation requires the detection
transport at that weight and the normalized-map comparison.
The formal binding here transports the displayed zero-level
infinite extension with the explicit comparison data.
A finite-quotient equation at any level is different: it first
requires a lift preserving both specified detections before this
full-object transport can be applied.
% Blueprint note: The original article number is 5.14; other numbering is generated by the blueprint. Throughout the article, spectra are assumed to be 2-completed, connective, and of finite type. Read this statement together with its original context.
% Blueprint note: \href{paper/main.tex.html#L1581}{Original statement, hypotheses, and context} · \href{paper/2412.10879.pdf}{Article PDF}
% Blueprint note: The printed boundary group on $Y$ has index $(s,t)$ although $y$ has degree $(s+n,t+n)$. The typed conclusion uses the latter degree and keeps zero and inessential extensions. See the statement-typing correction record.
% Blueprint note: The AimPaper statement is retained verbatim and directly binds its proved representative-form transport. The realization coordinates and actual comparison remain explicit inputs; the binding does not assert that the Standard comparison producer has been constructed.
% Blueprint note: The deferred completed-background migration is tracked in [issue 31](https://github.com/SII-MATH/Kervaire126Main/issues/31).
\lean{KIP126.ProofTools.Paper.infiniteExtension_to_classicalExtension}
\end{proposition}
\begin{proof}
\uses{sem:ProofTools:extension-representatives,sem:Literature:lambda-inversion-detection,sem:Literature:lambda-inversion-adams-comparison,sem:AdamsSynthetic:normalized-map,paper:3.11,paper:2.5}
\textit{(See \cite{LWX126}, Proposition 5.14.)} \href{paper/main.tex.html#L1589}{Original proof}.

Invert $\lambda$ in the filtered synthetic map $\widehat f$. The comparison with the classical map $f$ induces a morphism of extension spectral sequences. On the starting page, Proposition~\ref{paper:3.11} sends a label in $Z_\infty/B_{1+t-w}$ to its class in $Z_\infty/B_\infty=E_\infty$.

For the formal representative transport, put $e=e(f)$ and $w=t+e$.
The actual extension equation gives one representative
$\beta\in\pi_{t-s,w}(\Sigma^{0,e}\nu X)$ whose image under $\hat f$
has the specified target detection. Apply the detection transport theorem
to this same $\beta$ and to $\hat f_*\beta$. The prescribed $E_2$
comparisons return precisely $x$ and $y$.

Realization of the normalized factorization, followed by cancellation of
the invertible realized lambda map, gives
\[
\theta_Y\circ\operatorname{Re}(\hat f)
 =f\circ\theta_{\Sigma^{0,e}\nu X}.
\]
Thus the first realized representative is $a$, and the second is exactly
$f_*a$. The sphere-coordinate transport accounts for
$(t+n)-(s+n)=t-s$. By Proposition~\ref{paper:2.5}, these detections give
the classical $f$-extension in the conclusion. The target is taken in its
own bidegree $(s+n,t+n)$, as required for the quotient by its boundary
group. Localization can enlarge the boundary indeterminacy, so the
resulting extension need not be essential. The argument includes zero
and boundary labels and imposes no nonvanishing or essentiality condition.
\end{proof}

As observed in \cite{LWX126}, Remark 5.15, the classical comparison
in Proposition~\ref{paper:5.14} forgets the original lambda layer,
the finer synthetic target coset, and the chosen representative and
its image. Recovering these data requires a target-preserving lift.
The following criterion describes the corresponding additive
obstruction.

\begin{proposition}[Lifting a specified family of finite representatives]
\label{sem:ProofTools:ml-representative-family-lift}
\level{coarse}
An additive exactness criterion for the extra lifting step needed by
classical extensions and permanence outputs.\par
Let $q:D\to A$ and $\delta:A\to C$ be actual homomorphisms of abelian
groups with $\ker\delta=\operatorname{im}q$. Let $V\leq A$ be the
subgroup of permitted corrections preserving the specified source and
image labels, and fix $a_0\in A$. Then
\[
 (a_0+V)\cap\operatorname{im}q\ne\varnothing
 \quad\Longleftrightarrow\quad \delta(a_0)\in\delta(V).
\]
If $q$ is restriction from the full synthetic homotopy group, this
constructs a full lift of a corrected representative in the prescribed
family. If $q$ is only a restriction between two finite quotients, it
constructs only that one finite lift. Neither interpretation assumes
permanence of the output label.
\uses{sem:AdamsSynthetic:lambda-quotients}
\lean{KIP126.ProofTools.Standard.representativeFamilyLift_iff}
% Correspondence: The public theorem proves the displayed equivalence for the actual homomorphisms and correction subgroup. It constructs a lift of some member of the specified family, not necessarily of the initially chosen representative.
\end{proposition}
\begin{proof}
If $q(d)=a_0+v$, then $0=\delta(a_0)+\delta(v)$, so
$\delta(a_0)=\delta(-v)\in\delta(V)$. Conversely, write
$\delta(a_0)=\delta(v)$ with $v\in V$. Then
$\delta(a_0-v)=0$, and exactness gives $q(d)=a_0-v$ for some $d$.
Since $-v\in V$, this is a lift of a representative in the specified family.
\end{proof}

\section{Generalized Leibniz with actual level conditions}

\begin{theorem}[Generalized Leibniz with an arbitrary lambda level]
\label{sem:ProofTools:generalized-leibniz-at-level}
\level{coarse}
A level-indexed revision of \cite{LWX126}, Theorem 6.1.\par

The proof follows the quotient-boundary square argument of that theorem, with crossing hypotheses imposed on the actual level edges. In the fixed model, let $f:X\to Y$, $e=e(f)$, and
\[2\le n\le r,\qquad e\le m\le n-2+e,\qquad l\ge e.\]
Let $k\in\mathbb N$ be arbitrary, $q=n+k-1$, and $p=(t-s,t+e-k)$. Take
\[
x\in E_2^{s,t}(X),\quad y\in E_2^{s+m,t+m}(Y),\quad
x_\infty\in Z_\infty^{s+r,t+r-1}(X),\quad
y_\infty\in Z_\infty^{s+r+l,t+r+l-1}(Y).
\]
Suppose $d_r(x)=x_\infty$, there is the level-$k$ finite extension
\[d_m^{F_q}(\lambda^k x)=\lambda^{k+m-e}y,\]
and the existing zero-level infinite extension $d_l^F(x_\infty)=\lambda^{l-e}y_\infty$. Suppose also
\[
\bigl(NC(\delta_q^A;p,s,r)\lor NC(F_q;p,s,m)\bigr)
\land NC(\Sigma^{1,-q}F;p,s+r,l).
\]
Then the classical Adams differential is
\[d_{r+l-m}(y)=y_\infty.\]
There is no upper bound on $k$. The left edge and the propagated bottom edge are constructed in the proof; they are not extra caller premises.

\textit{Zero-level specialization and the original rule.}
At $k=0$, the finite equation is in quotient $n-1$. The preceding
representative comparison supplies this equation from an original
$(f,E_n)$-extension. The actual no-crossing hypotheses above must still
be verified on the indicated maps and weight. The original Theorem 6.1
instead uses classical crossing predicates; the unrestricted comparison
in \cite{LWX126}, Proposition 5.10, cannot replace those verifications.
Thus the zero-level specialization is a corrected rule with actual
crossing conditions, not a proof of the original unrestricted statement.
The original statement and its counterexamples are retained in the
article and the AimPaper corrections.
\uses{sem:ProofTools:level-finite-extension,paper:2.12,paper:4.6}
\lean{KIP126.ProofTools.Standard.generalizedLeibnizAtLevel}
% Blueprint note: This is the authorized level-indexed revision. It neither restores the retired original GL theorem nor identifies the old crossing predicates with the actual conditions. Its proof inherits the registered upstream construction and comparison obligations.
\end{theorem}
\begin{proof}
\textit{Proof sketch in the fixed model.}
Form the actual quotient-boundary square at $q=n+k-1$ and weight $t+e-k$. The classical differential constructs its left extension. Propagate the given infinite extension to the bottom edge at level $r-n$. These constructions retain the specified source and target labels. Apply the filtered-square extension theorem using the assumed no-crossing alternative on the left or top edge and the assumed actual bottom no-crossing. The resulting boundary equation, read through the weighted boundary comparison, is the classical differential $d_{r+l-m}(y)=y_\infty$. All steps occur at the displayed levels, so no upper bound on $k$ or cancellation of $\lambda$ is used.

More explicitly, the left equation has source $\lambda^kx$ and
target $\lambda^{r-n}x_\infty$ under the actual boundary shift.
The top equation is the displayed level-$k$ finite extension.
The infinite input propagates to the bottom by $\lambda^{r-n}$;
its source and target are transported into the shifted bottom map,
rather than by adding $k$ independently to every exponent.
Setting $k=0$ gives the same argument at zero level. It makes no
comparison with the classical crossing candidates of the original rule.
The finite-layer surjectivity composition below uses a finite
naturality square instead of assuming this rule's infinite input.
\end{proof}


\begin{remark}[A cutoff version of the Leibniz square]
\label{paper:6.5}
\level{coarse}
From \cite{LWX126}, Remark 6.5 and Theorem 2.12.\par
In the quotient-boundary square used for the Leibniz argument, let $x,y,z,w$ be the four specified infinite-page labels at a common stem and weight, with top, left and bottom extension lengths $m,r,l$. The four-edge theorem permits a cutoff $c$ with
\[
0<c\le r+l-m.
\]
Besides the three input extensions and no crossing on the top or left edge, it suffices to require that the bottom extension has no crossing hitting filtration at least $s+m+c$, and that the right edge has an extension of length $c-1$ from $y$ to zero with no crossing. Then the right edge has an extension of length $r+l-m$ from $y$ to $w$.

This is the extension-square generalization proposed in the remark. Reading its right edge as a specified classical Adams differential additionally requires the boundary and label comparisons for that square; the cutoff does not by itself supply them. At positive lambda level the
same four-edge statement must be instantiated at the actual quotient
and weight of the level square, including the right-edge zero extension.
At level zero it remains this cutoff specialization; no arbitrary-level
classical differential theorem is asserted solely by choosing $c$.
\uses{paper:2.12}
\lean{KIP126.ProofTools.Standard.cutoffLeibnizSquare}
% Correspondence: The complete cutoff square is publicly bound, including the actual quotient boundary, fixed weight and essential-crossing threshold.
\end{remark}
\begin{proof}
Apply Theorem~\ref{paper:2.12} to the actual quotient-boundary square, with its top length $m$, left length $r$ and bottom length $l$. Its right-edge zero extension raises the common composite to filtration $s+m+c$; the bottom no-crossing condition in that range identifies the composite's target with $w$. This gives the asserted right-edge extension. This argument retains the right-edge zero-extension hypothesis, rather than merely dropping part of the bottom no-crossing condition.
\end{proof}

\subsection{Why the crossing hypothesis is needed}

The following original example and remarks explain the crossing
hypothesis in the Leibniz rule and the role of the target coset.

\begin{proposition}[A crossing obstructing a Leibniz conclusion]
\label{paper:6.8}
\label{paper:6.6}
\label{paper:6.9}
\label{paper:6.10}
\label{exam:Leibnizno}
\label{rem:chuaerror}
\level{coarse}
From \cite{LWX126}, Example 6.8 and Remarks 6.6, 6.9--6.10.\par
Let $f=2:S^0\to S^0$, normalized by $[h_0]$. At
\[
n=r=2,\quad m=1,\quad l=2,\quad e(f)=1,\quad s=1,\quad t=16,
\]
the input relations are
\[
d_2(h_4)=h_0h_3^2,\qquad
d_1^{f,E_2}(h_4)=h_0h_4,\qquad
d_2^{f,E_\infty}(h_0h_3^2)=0.
\]
The bottom relation has the essential crossing
\[
d_1^{f,E_\infty}(d_0)=h_0d_0.
\]
Consequently, omitting the bottom no-crossing hypothesis from the original Leibniz rule would give the incorrect conclusion $d_3(h_0h_4)=0$.

The same shorter extension also makes
\[
d_2^{f,E_\infty}(h_0h_3^2)=h_0d_0
\]
inessential. Applying a rule without the required crossing control to these two choices would give the conflicting outputs
\[
d_3(h_0h_4)=0,\qquad d_3(h_0h_4)=h_0d_0.
\]
The maximal-image remark makes a stronger assertion about actual
homotopy classes. It chooses a lift of $h_0h_3^2$ in
$\pi_{14,17}S^{0,0}$ whose $[h_0]$-image is zero in
$\pi_{14,18}S^{0,0}$. A second lift has image $\lambda\beta$,
where $\beta\in\pi_{14,19}S^{0,0}$ has label $h_0d_0$.
These weights make explicit the lambda factor implicit in the
article's bracket notation. The zero-image choice is maximal in
the sense discussed there, yet the proposed rule would give the
incorrect zero conclusion. A maximal-image choice is not a
replacement for the crossing hypothesis. This actual zero-image
assertion is stronger than the length-two extension having target
label zero.
\uses{paper:5.4,paper:5.7,paper:5.8,paper:5.13}
\lean{KIP126.ProofTools.Standard.paper6_8_complete_crossing_warning}
% Correspondence: The public theorem contains the complete numerical warning: the actual normalized map, finite and infinite relations, both other no-crossing checks, the weight-18 essential crossing, inessentiality with nonzero infinity target, and two actual sphere lifts with respectively nonzero and zero h0-image. The same zero image has every lambda division. It uses the independent classical d3 input and does not invoke a discarded Leibniz rule or the disputed general crossing equivalence. The Chua locators remain those discussed in AimPaper.
\end{proposition}
\begin{proof}
\uses{sem:Literature:one-line-differentials,sem:Literature:two-line-classical-inputs,sem:Computation:bounded-comparison,sem:Literature:finite-bockstein,sem:Literature:infinite-nu-window,sem:Literature:selected-lifts,sem:Literature:selected-realization,sem:Literature:filtered-product-detection,sem:AdamsSynthetic:finite-separation}
\textit{The nonzero classical conclusion.} The independent classical input in
\ref{sem:Literature:two-line-classical-inputs} gives a nonzero $d_3$ from
the specified $h_0h_4$ to degree $(5,19)$. The bounded Ext archive gives
one basis element there, namely $h_0d_0$, with $d_0$ the fixed archived
generator in degree $(4,18)$. A one-dimensional $\mathbb F_2$ space has
exactly one nonzero element, so the target is $h_0d_0$. In particular the
zero conclusion is false. This proves that conclusion without assuming
the extension relations whose use is being examined.

\textit{The extension representatives.} For the actual normalized
map, the first finite relation comes from multiplication by $h_0$
on the first lambda quotient. Apply the specified Bockstein
formulas to the independent differentials
$d_2(h_4)=h_0h_3^2$ and $d_3(h_0h_4)=h_0d_0$.
Their naturality constructs full classes $a$ and $\beta$ with
the prescribed labels and $[h_0]a=\lambda\beta$.
The class $d_0$ has a full lift $D$: all its outgoing differentials
vanish, by the bounded Ext calculation and the sphere vanishing line.
The actual first-quotient product shows that $[h_0]D$ has label
$h_0d_0$. Hence $a-\lambda D$ keeps the original source label,
while its image is
\[
\lambda\bigl(\beta-[h_0]D\bigr)\in F^6\pi_{14,18}S^{0,0}.
\]
It therefore has zero target label in filtration five. This
constructs the zero-target relation; it does not assert that this
actual image is zero. The length-one extension from $d_0$ has
nonzero target and is essential in its minimal-length range.

\textit{Crossing and inessentiality.} Express the shorter relation
at weight $18$ using $\lambda D$ and $\lambda[h_0]D$.
Its target grade is nonzero because $h_0d_0\notin B_2$, as follows
from the nonzero classical $d_3$. In the minimal-length range,
the actual no-crossing criterion proves this relation essential.
This is a fresh check at weight $18$, not an assertion that
multiplication by lambda always preserves essentiality.
It is the required crossing. The two long relations have the
same source infinite-page class, since their actual sources
differ in higher filtration. Uniqueness of the extension
differential, compared with the zero-target witness, shows that
every long witness with target $h_0d_0$ is inessential.
Its target infinite-page class is nevertheless nonzero.

\textit{The actual zero-image lift.} The boundary construction
of $a$ gives $\lambda a=0$, so its realization is zero.
The class $d_0$ survives nontrivially. Product detection and the
fact that $h_0d_0$ is a classical boundary show that the
realization of $2D$ is infinitely Adams-filtered. Ordinary
sphere separation makes it odd torsion; choose an odd integer
$N$ annihilating it. Replacing $D$ by $ND$ preserves its
first-quotient label. Put
\[
z=[h_0](a-\lambda ND)\in F^6\pi_{14,18}S^{0,0}.
\]
This class has zero realization. In this stem the filtration-six
Ext group is generated by $h_0^2d_0$, a $d_3$ target, and all
higher filtration groups vanish. The corresponding synthetic
window therefore puts $z$ in every filtration. Zero realization
also makes $z$ lambda-torsion.

For weights at most $18$, the actual first-quotient exact sequence
shows that the lambda kernel in this stem has no infinitely
filtered nonzero class. At weights $17$ and $18$, its possible
generators are the boundaries of $h_4$ and $h_0h_4$;
their specified nonzero associated grades exclude infinite
filtration. At lower weights the relevant first-quotient group
vanishes. Induction on a lambda-annihilating power, which only
lowers weight, now gives $z=0$. Thus $a-\lambda ND$ has the
original source label and actual zero $[h_0]$-image. The fixed
sphere and suspension comparisons transport this equation to
the displayed $S^{0,0}$ coordinates.
That same zero image is divisible by every power of lambda,
so it is a maximal-image choice. The nonzero classical $d_3$
already proved above still excludes the proposed zero conclusion.

This argument concerns the specified zero-layer example. Transporting its relations to another quotient or weight does not by itself preserve essentiality or no-crossing, so it is not a claim of the same counterexample at every retained level.
\end{proof}

\section{The finite-layer generalized Mahowald trick}

The layer parameters remain visible in the inputs and the output.
The external May input is instantiated in the actual quotient diagram;
it is not replaced by a second, level-indexed May theorem.

% AimPaper: 6.12 theorem 1780 1799
\begin{theorem}[The finite-layer generalized Mahowald trick]
\label{paper:6.12}
\source{A finite-layer revision of AimPaper, Theorem 6.12}
\level{coarse}
A finite-layer revision of \cite{LWX126}, Theorem 6.12.\par
Fix one classical triangle $X\xrightarrow fY\xrightarrow gZ\xrightarrow h\Sigma X$ and its compatible normalized synthetic triangle. Write $Q_DW=\nu W/\lambda^D$, retaining the fixed normalization and suspension comparisons in every map. Put
\[
\begin{gathered}
n_1=n-e_f\geq1,\quad m_1=m-e_g\geq0,\\
k=l-e_h\geq0,\quad r=n+m+l,\qquad e_f+e_g+e_h=1.
\end{gathered}
\]
For $i,j\geq0$ with $j<n_1$, set
\[
\begin{gathered}
A=n_1+k+i-j,\quad B=m_1+j+1,\quad N=A+B=r+i,\\
K=i+k,\quad J=A+j,\quad Q=N-K=r-k.
\end{gathered}
\]
Thus $0\leq i,K<A$, $0\leq j<B$, and $0\leq J<N$. Take classical labels in the prescribed bidegrees
\[
\begin{gathered}
\bar x\in Z_{r-1}^{s,t}(Z),\quad \bar y\in Z_\infty^{s+r,t+r-1}(Z),\\
x\in Z_{n_1-j}^{s+l,t+l-1}(X),\quad y\in Z_{m_1+1}^{s+l+n,t+l+n-1}(Y),
\end{gathered}
\]
with $d_r\bar x=\bar y$. Write the normalized vertices as $U=\Sigma^{0,e_f}\nu X$, $V=\nu Y$, and $W=\Sigma^{0,-e_g}\nu Z$. Choose actual suspension--quotient isomorphisms
\[
\varepsilon_X^D:Q_D(\Sigma U)\cong\Sigma Q_DU\quad(D=A,N),\qquad
\varepsilon_Y^N:Q_N(\Sigma V)\cong\Sigma Q_NV.
\]
Each must commute with its quotient inclusion. Compatibility with reduction $Q_N\to Q_A$ on bigraded sphere representatives follows internally from these inclusion squares and the finite no-return comparison; it is not an extra hypothesis. The assertion below is quantified over these choices. The map $h_A$ uses $\varepsilon_X^A$, and the output uses $\varepsilon_X^N,\varepsilon_Y^N$; a comparison preserving only objects is insufficient. Finite source labels are transported by these actual shifts and quotients, using the intrinsic finite source label rather than an independently chosen positive-shift $E_2$ identification.

Their actual finite infinite-page classes must support the following three equations in the indicated fixed maps:
\[
\begin{gathered}
d_l^{h_A}(\lambda^i\bar x)=\lambda^Kx,\\
d_r^{\delta_{A,B}}(\lambda^i\bar x)=\lambda^{B-1}\bar y,\\
d_m^{g_B}(\lambda^jy)=\lambda^{B-1}\bar y.
\end{gathered}
\]
The first two equations use the same infinite-page source class. The last two use the same top-layer target class. The $h$ target is related to the specified $x$ by the fixed actual suspension comparison. Assume actual no-crossing on either $h_A$ or $\delta_{A,B}$: the entire triangle $s<u\leq v\leq s+l$ or $s<u\leq v\leq s+r$, respectively. Both are at the same actual stem and weight. In unshifted $Z$ coordinates this is $(t-s,t-i)$; in the normalized third vertex it is $(t-s,t-i-e_g)$.

Then there exist the same actual representatives $a,b,c$, with
\[
\begin{gathered}
a\text{ detecting }\lambda^i\bar x,\\
h_A(a)\text{ detecting }\lambda^Kx,\\
\delta_{A,B}(a)=g_B(b)\text{ detecting }\lambda^{B-1}\bar y,\\
b\text{ detecting }\lambda^jy,
\end{gathered}
\]
\[
\rho(c)=h_A(a),\qquad f_N(c)=-\lambda^Ab,
\]
such that $c$ detects $\lambda^Kx$ and its actual $f_N$ image detects $\lambda^Jy$. Equivalently for the resulting page relation,
\[
x\in Z_Q,\qquad d_n^{\hat f_N}(\lambda^Kx)=[\lambda^Jy].
\]
The source label lies in $Z_Q/B_{K+1}$; the target lies in $Z_{m_1+1}/B_{J+1}$ modulo shorter actual $\hat f_N$ extension images. The actual representatives are retained in addition to this quotient relation. At $i=j=0$ the output is in the strong quotient $Q_r$:
\[
\begin{gathered}
x\in Z_{n+m+e_h},\\
d_n^{\hat f_r}(\lambda^kx)=[\lambda^{r'-1}y],\quad r'=n_1+k+1.
\end{gathered}
\]
A zero target label means a higher filtration, not a zero homotopy image. No power of $\lambda$ is cancelled.

\textit{Relation with the zero-level statement.}
Setting $i=j=0$ leaves the output lambda level $k=l-e_h$.
Only when $k=0$ as well does this output lie at level zero.
For $k>0$, the quotient has length $r$ and the source denominator
is $B_{k+1}$; the original statement would require a specified
zero-level source in the smaller quotient of length $r-k$.
It also requires the prescribed target modulo $B_{r'}$ there.
Neither assertion follows by cancelling $\lambda^k$.
The conditional division theorem retains actual source and image
equations under its extra kernel condition, but does not select
arbitrary low-layer source and target labels.
The original formula and the failure of its unrestricted division
are recorded in the article and the AimPaper corrections.
\uses{sem:ProofTools:mh-actual-no-crossing,sem:ProofTools:sphere-label-detection,sem:ProofTools:mh-common-source,sem:Literature:normalized-finite-quotient-row,paper:3.20,paper:6.11}
\sketch
\lean{KIP126.ProofTools.Standard.generalizedMahowaldAtLevel}
% Blueprint note: The complete finite-layer assertion is proved using the fixed Literature row, without a caller-supplied horizontal-row premise. It retains actual no-crossing synchronization, the same common representative, both signed image equations, both exact output detections, the original source cycle qualification and the specified target strong quotient. The proof introduces no local gap or axiom, does not use the frozen old MH or GL law, and still inherits the registered upstream model and comparison obligations.
% Blueprint note: This replaces the unprovable zero-layer version formerly displayed here. The original Theorem 6.12 remains in AimPaper and in the corrections page. The historical Lean law is not used as a proof of this revised target. A classical extension would additionally require an appropriate target-preserving full lift; source permanence alone does not provide it.
\end{theorem}
\begin{proof}
\textit{(See \cite{LWX126}, Theorem 6.12.)} \href{paper/main.tex.html#L1860}{Original Mahowald argument (before the finite-layer correction)}.

\textit{Proof of the finite-layer revision.} Fix the compatible normalized triangle and the data of the theorem. Actual no-crossing gives one representative $a$ for which the prescribed $h_A$ image and the $\delta_{A,B}$ image have the required detections. This synchronization is the common-source construction at \ref{sem:ProofTools:mh-common-source}; it does not replace either specified label.

Top-layer uniqueness identifies the boundary image of this $a$ with $g_B(b)$ for the chosen representative $b$. Apply the fixed finite-quotient horizontal row of \ref{sem:Literature:normalized-finite-quotient-row} and the triangle construction of Lemma~\ref{paper:6.11}. They produce a single $c$ satisfying both prescribed signed image equations. The comparison of the cofiber objects preserves both structure arrows and the same suspension comparison, so these are equations in the required quotients.

The normalized no-return argument removes the restriction defect of $c$. Visible-layer quotient reflection and the source/image detection comparisons then identify its two output detectors and retain the cycle qualification of the original source. Reading the second image equation in the specified strong quotient gives the finite extension in the conclusion. No $\lambda$ layer is cancelled, and the target is not projected to a smaller quotient. This proves the revised assertion relative to its registered upstream model and comparison inputs; it does not prove the zero-layer statement printed in the original Theorem 6.12.
\end{proof}

\subsection{Cofiber examples of the Mahowald rule}

\begin{proposition}[The nu-cofiber calculation of a finite extension]
\label{paper:6.19}
\label{exam:Mahowald}
\level{coarse}
From \cite{LWX126}, Example 6.19.\par
Consider
\[
S^3\xrightarrow{\nu}S^0\xrightarrow{i}C\nu\xrightarrow{q}S^4,
\qquad (e(\nu),e(i),e(q))=(1,0,0).
\]
In Atiyah--Hirzebruch notation, $a[4]$ has top-cell image $a$ and $b[0]$ denotes the bottom-cell image of $b$. The input data are
\[
\begin{aligned}
d_0^{q,E_2}(h_0h_4^2[4])&=h_0h_4^2,\\
d_2(h_0h_4^2[4])&=h_0p[0],\\
d_0^{i,E_2}(h_0p)&=h_0p[0],
\end{aligned}
\]
where $h_0p[0]$ is a permanent cycle. The classical $d_2$ and the length-zero $q$-extension have no crossing in the sense used in the article. With
\[
n=2,\quad m=l=0,\quad n_1=1,\quad m_1=l_1=0,
\quad r=r'=2,
\]
the claimed output is that $h_0h_4^2\in E_2^{3,36}(S^3)$ survives to $E_3$ and
\[
d_2^{\nu,E_3}(h_0h_4^2)=h_0p,
\qquad h_0p\in E_2^{5,38}(S^0).
\]
The qualification on $E_3$ allows the zero page class: the specified
$E_2$ label $h_0h_4^2$ is itself a $d_2$ boundary. The finite synthetic
extension retains that label; no nonzero classical $E_3$ survival is asserted.
\uses{paper:5.4,sem:Literature:nu-cofiber,paper:4.14}
\lean{KIP126.ProofTools.Standard.paper6_19_finite_nu_extension,KIP126.ProofTools.Standard.paper6_19_actual_cofiber_d2}
% Correspondence: The two public bindings separately prove the finite extension with its E3 cycle qualification, and the existence of an actual cofiber class with the displayed top-cell image, cofiber d2 and permanent bottom-cell target. Both use the same independently constructed full representatives. They do not identify an arbitrary preselected AH lift, assert a nonzero bottom-cell target, or reproduce the article's secondary-Steenrod-algebra calculation. This n=2 result is distinct from the existing n=3 Main theorem.
\end{proposition}
\begin{proof}
\uses{paper:6.24,paper:6.11,paper:6.12,sem:Literature:finite-bockstein,sem:ProofTools:extension-representatives}
\textit{An independent proof of the finite output.} The direct Bockstein
construction in Example~\ref{paper:6.24} gives one full source $a$ and
one full target representative $\beta$, with the prescribed labels and
$\widehat\nu(a)=\lambda\beta$. Project these same representatives to
$Q_2$. Naturality preserves their image equation and their finite
detectors. Since $e(\nu)=1$, the length bound on effective page $3$ is
$2\le3-2+1$, with equality. This is exactly the finite extension
displayed above, in the strong quotient of length $2$. The source
label is a permanent cycle and hence reaches $E_3$, with zero page
class allowed. This proof uses the independent sphere $d_3$ input;
it does not derive that input by the original Leibniz argument.

\textit{The cofiber data.} They can also be recovered from the same
full representatives, without assuming the cofiber differential.
Rotate the fixed normalized $\nu$ triangle and pair it with the
actual $\lambda$-cofiber triangle. The common-lift construction
gives one class in $Q_1C\nu$ with the prescribed suspended top-cell
image and boundary equal to the bottom-cell image of the chosen
full target. Its restriction defect vanishes by the finite
no-return comparison. The specified-lift Bockstein formula then
gives the displayed $d_2$ and proves that its bottom-cell target is
permanent. This constructs an appropriate Atiyah--Hirzebruch lift;
it does not identify every lift with the same top-cell label.

The article obtains these data by a secondary-Steenrod-algebra
calculation. The argument here instead uses the independent sphere
$d_3$ input already used for the full representatives. The original
page-two crossing checks are vacuous. In the finite-layer Mahowald
rule, both input levels and the output level are zero, giving the
same quotient length $2$. No unresolved general crossing
equivalence is an input to either construction above.

For positive input level $a$, the corresponding finite target relation
is $\mathcal E^\nu_{2+a,a}(2;h_0h_4^2,h_0p)$.
Applying the level-indexed Mahowald theorem requires the associated
actual input representatives and their comparisons; the zero-level
calculation alone is not their formal producer.
The separate theorem below has length three and different bidegrees.
It does not establish this length-two numerical example.
\end{proof}

\begin{theorem}[The zero-layer nu-cofiber Mahowald instance]
\label{sem:ProofTools:nu-zero-layer}
\level{coarse}
In the fixed cofiber triangle for $\nu$, take labels $x\in E_2^{8,133}(S^3)$, $y\in E_2^{11,136}(S^0)$, $\bar x\in E_2^{8,134}(C\nu)$ and permanent $\bar y\in E_2^{11,136}(C\nu)$. Require $x\in Z_2$, $y\in Z_1$, $\bar x\in Z_2$, the actual suspension label for $x$, the length-zero connecting extension on page $3$, $d_3(\bar x)=\bar y$, its specified no-crossing alternative, and the length-zero bottom-cell extension on page $2$. Then $x\in Z_3$ and there is an actual $\nu$-extension of length $3$ on page $4$ from $x$ to some $y'$ with $y'-y\in B_3$.
\uses{sem:Literature:nu-cofiber,sem:Literature:connected-homology,sem:AdamsSynthetic:shift-coherence,sem:AdamsSynthetic:ess-shifts,sem:ProofTools:common-labels,sem:ProofTools:filtration-representatives,sem:ProofTools:extension-representatives,sem:ProofTools:sphere-label-detection,sem:ProofTools:finite-strip,paper:5.4}
\lean{KIP126.ProofTools.Internal.Mahowald.main_nu_mahowald_zero_layer}
% Blueprint note: The instance has its own proof. Its fixed-model constructions and comparison inputs retain their registered upstream obligations.
\end{theorem}
\begin{proof}
\textit{Proof sketch in the fixed model.}
Use the actual quotient-boundary diagram of the fixed $\nu$ cofiber, with its specified suspension comparison. The classical $d_3$ supplies the boundary equation, and the two length-zero extensions supply the connecting and bottom-cell equations. The no-crossing alternative synchronizes the required representatives. In this particular finite diagram, the quotient-window bounds identify the source as a $Z_3$ label and identify the resulting length-three extension target modulo $B_3$. Thus the extension is on page $4$ with the stated congruence. This is an independent argument for this numerical instance, not cancellation of a $\lambda$ power in the general Mahowald theorem.
\end{proof}

\section{Restriction, indeterminacy and finite stretching}

Restriction, changes in essentiality, and stretching are separate
questions. The following results use zero-level page extensions with
their stated domains. Positive-level versions require the corresponding
label comparisons; none of these results is an automatic transport
of all crossing information.

\begin{theorem}[Descent of finite-page extensions in the valid length range]
\label{sem:ProofTools:finite-extension-descent}
\level{coarse}
A page-range revision of \cite{LWX126}, Proposition 6.20, first assertion, with the page--length domain of Definition 5.4.\par

In the fixed model, let $X,Y$ be classical objects and $f:X\to Y$ a map, and put $e=e(f)$. Take finite pages $r_{\rm early},r_{\rm late}\in\mathbb N$, integers $n,s,t$, and common labels
\[
x\in E_2^{s,t}(X),\qquad y\in E_2^{s+n,t+n}(Y).
\]
Suppose that
\[
2\leq r_{\rm early}\leq r_{\rm late},\qquad
n\leq r_{\rm early}-2+e,
\]
and that there is an extension, in the sense of Definition 5.4,
\[
d_n^{f,E_{r_{\rm late}}}(x)=y.
\]
Then the same map and the same specified labels admit the earlier-page extension
\[
d_n^{f,E_{r_{\rm early}}}(x)=y.
\]
The late extension includes the lower length bound $e\leq n$ and the cycle conditions
\[
x\in Z_{r_{\rm late}-1}^{s,t}(X),\qquad
y\in Z_{r_{\rm late}-1-n+e}^{s+n,t+n}(Y).
\]
The conclusion includes the corresponding earlier-page cycle conditions. The extension need not be essential; no crossing hypothesis is imposed.

\textit{Level comparison.}
This assertion uses zero-level page extensions. At fixed positive
level $k$, the corresponding quotients would have lengths
$r_{\rm late}+k-1$ and $r_{\rm early}+k-1$, at the same weight
$t+e-k$. A positive-level restriction must compare the prescribed
source and target labels along that actual reduction. The zero-level
result does not assert this additional representative comparison,
or preservation of essentiality or no-crossing.
\uses{paper:5.4}
\lean{KIP126.ProofTools.Paper.finiteExtension_descends}
% Blueprint note: This corrects the unrestricted first assertion of Proposition 6.20. It imposes no essentiality or crossing premise and retains the registered upstream obligations.
\end{theorem}
\begin{proof}
Apply quotient reduction to the source and target of the late normalized extension equation. Naturality gives the equation in the earlier quotient. The explicit upper bound $n\le r_{\rm early}-2+e$ places it in the valid length window of Definition~\ref{paper:5.4}; the late cycle conditions supply the corresponding early cycles. Reading the reduced equation gives the stated early extension, with the same map and labels.
\end{proof}



% AimPaper: 6.20 proposition 1981 1993
\begin{proposition}[Fixed-target indeterminacy after restriction]
\label{paper:6.20}
\source{A fixed-target revision of AimPaper Proposition 6.20, second assertion}
\level{coarse}
A fixed-target revision of \cite{LWX126}, Proposition 6.20, second assertion, with the late-cycle exclusion removed.\par
\label{prop:dec738d3}
Work in the fixed Standard model. Let $f:X\to Y$, put $e=e(f)$, and take finite pages $r',r\in\mathbb N$ and integers $n,s,t$ with
\[
2\le r'\le r,\qquad e\le n\le r'-2+e.
\]
Fix common labels
\[
x\in Z_{r-1}^{s,t}(X),\qquad y\in Z_{r-1-n+e}^{s+n,t+n}(Y),
\]
and witnesses $W_r,W_{r'}$ for $d_n^{f,E_r}(x)=y$ and
\begin{equation}\label{eq:dr'}
d_n^{f,E_{r'}}(x)=y.
\end{equation}
Both witnesses include the full cycle and length qualifications of Definition~\ref{paper:5.4}; the earlier relation is supplied by \ref{sem:ProofTools:finite-extension-descent}. Suppose that $W_r$ is essential.

Put $q=r'-1$, $A=\Sigma^{0,e}\nu X$, $B=\nu Y$, and let $F_q:A/\lambda^q\to B/\lambda^q$ be the actual normalized map. Retain the target class
\[
\bar y_{r'}\in E_\infty^{s+n,t+n,t+e}(B/\lambda^q)
\]
of the chosen witness $W_{r'}$, represented by the specified label $\lambda^{n-e}y$. Then the extension in (\ref{eq:dr'}) is inessential if and only if there exist
\[
a\in\mathbb N,\qquad 0<a\le n-e,\qquad
x'\in Z_{r'-1-a}^{s+a,t+a}(X),
\]
and a class $z\in E_\infty^{s+a,t+a,t+e}(A/\lambda^q)$ represented by the actual correction label $\lambda^a x'$, transported by the shift comparison and quotient isomorphism into the normalized source quotient at weight $t+e$, such that in this same extension spectral sequence
\[
d_{n-a}^{F_q}(z)=\bar y_{r'}.
\]
This is the full extension relation, including its source qualification. The shorter extension may be inessential and its source class may be zero. No exclusion from the later cycle set is imposed; the target is the original specified target, with $b=0$.

\textit{Level comparison.}
At level $k$, the same question belongs in the early quotient
$q=r'+k-1$ and weight $t+e-k$, retaining its actual target class.
The general representative criterion for inessentiality still concerns
a higher-filtration source mapping to that target. Expressing it by
classical correction labels requires the level-$k$ windows, their
boundary denominators and the comparison with the later quotient.
The present theorem proves the displayed zero-level fixed-target
criterion. It does not prove the original general-$b$ equivalence
with the later-cycle exclusion, nor its arbitrary-level extension.
\uses{paper:5.4,sem:ProofTools:finite-extension-descent,sem:ProofTools:extension-representatives}
\lean{KIP126.ProofTools.Paper.finiteExtension_inessential_iff_shorter}
% Correspondence: complete for the authorized fixed-target revision of the second assertion.
% Blueprint note: The revised first assertion remains stated and bound at sem:ProofTools:finite-extension-descent. The separate difference-preserving alternative is sem:ProofTools:nonliftable-or-crossing. Neither result proves the original unrestricted set-difference equivalence with general b. The local proof inherits the registered upstream obligations.
% Blueprint note: \href{paper/main.tex.html#L1981}{Original Proposition 6.20 and its context} · \href{paper/2412.10879.pdf}{Article PDF}. The original text remains in AimPaper and in the corrections tab.
\end{proposition}
\begin{proof}
\textit{Proof of the fixed-target revision; compare \cite{LWX126}, Proposition 6.20.} \href{paper/main.tex.html#L1995}{Original argument}.

The page--length domain condition $n\le r'-2+e$ allows restriction to the earlier page. Essentiality of the late witness implies that the specified early target is nonzero. The representative criterion for inessentiality in the actual early ESS identifies an inessential extension with an extension to that same target from a source in higher filtration. The minimal-length bound gives $0<a\le n-e$. The finite-window comparison expresses this higher source by a classical label $x'\in Z_{r'-1-a}$ and its actual correction label, retaining the same quotient, stem and weight. This proves the forward implication, including the source qualification.

Conversely, the shorter extension puts the same target in the earlier ESS's shorter-extension indeterminacy, so $W_{r'}$ is inessential. No claim that $x'$ fails to lift to the later cycle set is needed. Imposing that exclusion requires the separate late-crossing alternative at \ref{sem:ProofTools:nonliftable-or-crossing}; the original general-$b$ equivalence remains outside this revision.
\end{proof}

\begin{theorem}[A nonliftable short extension or a late crossing]
\label{sem:ProofTools:nonliftable-or-crossing}
\level{coarse}
Consider finite witnesses for the same $f,x,y,n$ on pages $r_{\rm early}\leq r_{\rm late}$. If the late witness is essential and the early one is inessential, then either a shorter early extension retains the original target and has a source in the earlier cycle set but outside the later cycle set, or the same late normalized extension spectral sequence has a nonzero higher-source crossing in the specified range. The first branch uses the actual correction label with $b=0$ and all length and cycle bounds.
\uses{sem:ProofTools:extension-representatives,sem:ProofTools:finite-strip,paper:6.20}
\lean{KIP126.ProofTools.Paper.finiteExtension_nonliftableShort_or_lateCrossing}
% Blueprint note: This disjunction is distinct from the original general-b set-difference equivalence. Its correction-label branch has b=0 and the original specified target; registered upstream obligations remain.
\end{theorem}
\begin{proof}
\textit{Proof sketch in the fixed model.}
Use the fixed-target shorter-extension criterion in the early quotient. If its higher source has no corresponding later cycle label, it is the first alternative. Otherwise lift the source label and compare the actual image with the fixed late target. Their nonzero discrepancy has a first grade in the finite higher-source range and gives the second alternative, an actual late crossing. Thus imposing the cycle-set difference requires retaining the crossing branch.
\end{proof}

% AimPaper: 6.23 corollary 2074 2081
\begin{corollary}[Finite-page stretching with a same-quotient obstruction]
\label{paper:6.23}
\source{A finite-page sufficient-condition revision of AimPaper Corollary 6.23}
\level{coarse}
A finite-page sufficient-condition revision of \cite{LWX126}, Corollary 6.23.\par
\label{cor:dfc6043e}
Work in the fixed Standard model. Let $f:X\to Y$, put $e=e(f)$, and let $r',r\in\mathbb N$ and $n,s,t\in\mathbb Z$ satisfy
\[
2\le r'\le r,\qquad e\le n\le r'-2+e.
\]
Suppose that the specified labels satisfy the later-page cycle conditions
\[
x\in Z_{r-1}^{s,t}(X),\qquad
y\in Z_{r-1-n+e}^{s+n,t+n}(Y),
\]
and that there is an earlier-page extension $d_n^{f,E_{r'}}(x)=y$.

Let $A=\Sigma^{0,e}\nu X$, $B=\nu Y$, and $q=r'-1$. Assume that there is no nonliftable obstruction in the same early quotient. Explicitly, exclude every pair $a,b\in\mathbb N$ with
\[
0<a\le r'-2,\qquad a+b\le n,
\]
and classes
\[
\xi\in E_\infty^{s+a,t+a,t+e}(A/\lambda^q),\qquad
0\ne\eta\in E_\infty^{s+n-b,t+n-b,t+e}(B/\lambda^q),
\]
such that $\xi$ is outside the image of the actual quotient reduction
\[
\rho_\infty:E_\infty^{s+a,t+a,t+e}(A/\lambda^{r-1})
\longrightarrow E_\infty^{s+a,t+a,t+e}(A/\lambda^q)
\]
and, in the same earlier-quotient ESS at stem $t-s$ and weight $t+e$,
\[
d_{n-a-b}^{\hat f_q}(\xi)=\eta.
\]
These obstructions include inessential extensions, $b=0$ and length zero; source nonliftability refers to the actual $\rho_\infty$ image, not to a difference of classical cycle sets. Then the same labels admit the later-page extension
\[
d_n^{f,E_r}(x)=y.
\]
The conclusion asserts existence of a later extension; it does not prescribe a lift of an arbitrarily chosen strict representative or a compatible infinite tower.

\textit{Level comparison.}
For a fixed positive level $k$, the analogous finite problem uses
$q'=r'+k-1$, $q=r+k-1$ and weight $t+e-k$.
Its obstruction must still be stated using the actual map
$\rho_\infty$ and all the indicated nonzero shorter relations in
the same early quotient. Merely replacing every exponent in the
zero-level formula by an exponent plus $k$ does not establish the
label comparisons. No such arbitrary-level extension, and no
case with $r=\infty$, is supplied by the bound finite theorem.
\lean{KIP126.ProofTools.Standard.finitePageExtensionStretching}
% Correspondence: complete for the stated finite-page sufficient-condition revision.
% Blueprint note: This revision strengthens the printed hypotheses by requiring the later target cycle qualification and excluding all specified nonzero same-quotient obstructions. It does not assert the original essential-crossing-only criterion or the case r=infinity. Its local proof is independent of the revised Proposition 6.20 and inherits registered upstream obligations.
% Blueprint note: \href{paper/main.tex.html#L2074}{Original Corollary 6.23 and its context} · \href{paper/2412.10879.pdf}{Article PDF}. The original text remains in AimPaper and in the corrections tab.
\end{corollary}
\begin{proof}
\uses{sem:ProofTools:finite-correction,sem:ProofTools:finite-strip,sem:ProofTools:extension-representatives}
\textit{Proof of the finite-page revision; compare \cite{LWX126}, Corollary 6.23.} \href{paper/main.tex.html#L2074}{Original statement and context}.

Use actual representatives of the earlier extension. The prescribed late source and target cycle qualifications supply lifts of their infinite-page labels along quotient reduction. Correct the difference between the chosen early source and a late lift over the finite filtration interval. At each step, a residual source with nonzero target that cannot be lifted would give exactly one of the excluded same-quotient obstructions. Every other residual can be corrected to the next filtration or already has image above the prescribed target filtration. The finite correction theorem and the finite-quotient cutoff therefore supply a late representative with the specified source and target labels. Reading its actual map equation gives the required extension.

This is an independent representative proof of the finite sufficient condition, not a contrapositive of the revised Proposition~\ref{paper:6.20}. It proves the stated revision, not the full printed corollary with only essential crossings excluded or an infinite later page.
\end{proof}

\subsection{The original infinite-extension application}

The article obtains this infinite-page application by stretching.
The proof below constructs the full synthetic representatives directly
from the independent sphere differentials.

\begin{proposition}[The infinite-page nu-extension used in the Leibniz example]
\label{paper:6.24}
\label{exam:stretchext}
\level{coarse}
From \cite{LWX126}, Example 6.24.\par
For $f=\nu:S^3\to S^0$, the finite relation of Example~\ref{paper:6.19},
\[
d_2^{f,E_3}(h_0h_4^2)=h_0p,
\]
is asserted in the article to extend to
\[
d_2^{f,E_\infty}(h_0h_4^2)=h_0p.
\]
The latter means that representatives in the full synthetic spheres have the prescribed source and image detections; it is stronger than the equation in the quotient by $\lambda^2$ alone.
\uses{paper:5.4}
\lean{KIP126.ProofTools.Standard.paper6_24_infinite_nu_extension}
% Correspondence: The public theorem proves the complete InfiniteExtension statement from a full source and its actual normalized-nu image, with both specified first-quotient labels and a strict image equality. It uses the independent classical d3 input, not the original Leibniz derivation or an assumed infinite lift.
\end{proposition}
\begin{proof}
\uses{sem:Literature:one-line-differentials,sem:Literature:two-line-classical-inputs,sem:Computation:bounded-comparison,sem:Literature:finite-bockstein,sem:ProofTools:sphere-label-detection}
\textit{Direct construction of the full representatives.}
Write $A=\Sigma^{0,1}\nu S^3$, $B=\nu S^0$ and let
$F:A\to B$ be the actual normalized $\nu$ map. For an actual quotient
boundary write
\[
\partial_q:\pi_{a,w}(Q_qV)\longrightarrow\pi_{a-1,w+q}(V).
\]
The independent inputs are $d_2(h_5)=h_0h_4^2$ and
$d_3(h_2h_5)=h_0p$; the latter target is fixed by the one-dimensional
Ext comparison used in the independent proof of Example~\ref{paper:6.7}.
The adjacent Bockstein comparison gives classes
\[
\xi\in\pi_{34,36}(Q_1A),\quad
z\in\pi_{34,36}(Q_2B),\quad
a=\partial_1\xi\in\pi_{33,37}(A),\quad
\beta=\partial_2z\in\pi_{33,38}(B).
\]
Here $a$ has the prescribed $h_0h_4^2$ source detector, transported by
the actual threefold suspension and normalization comparisons, and
$\beta$ has detector $h_0p$.

The actual first-quotient map sends $\xi$ to the class $h_2h_5$,
so that $F_1\xi=\rho_{1,2}z$. This is a comparison of operations:
quotient multiplication and its unit identify the action with
multiplication by $h_2$; the one-dimensional first-quotient source
fixes the sphere automorphisms in the suspension comparisons.
It is stronger than merely identifying the normalized map on the unit.
Boundary naturality and compatibility with restriction now give
\[
F(a)=\partial_1(F_1\xi)
     =\partial_1(\rho_{1,2}z)
     =\lambda\,\partial_2z=\lambda\beta.
\]
Thus a single full source and its actual image have the required
detectors. The source has filtration $3$, and the image has the
specified filtration-$5$ detector $\lambda h_0p$ at weight $37$.
This is precisely the displayed length-two infinite-page extension.
The classical page classes may vanish; the specified $E_2$ labels
and their synthetic detectors are retained throughout.

\textit{Relation to the article's stretching argument.}
The article instead applies Corollary~\ref{paper:6.23} with later page
$\infty$, using Example~\ref{paper:5.13}. The revised finite-page
stretching theorem alone does not provide that passage. The direct
construction above proves this numerical full extension without
choosing unrelated lifts at successive finite levels, and does not
assert a general finite-to-infinite stretching theorem or an
arbitrary-level version of the example.
\end{proof}

\subsection{The original Leibniz application}

The following source calculation uses the extension in the preceding
example, so it is placed after that example rather than before its
input has been described. Its general rule is
\ref{sem:ProofTools:generalized-leibniz-at-level}.

\begin{proposition}[A classical d-three obtained by the Leibniz argument]
\label{paper:6.7}
\label{exam:Leibnizyes}
\level{coarse}
From \cite{LWX126}, Example 6.7.\par
For $f=\nu:S^3\to S^0$, consider
\[
\begin{gathered}
x=h_5\in E_2^{1,35}(S^3),\qquad
y=h_2h_5\in E_2^{2,36}(S^0),\\
x_\infty=h_0h_4^2\in Z_\infty^{3,36}(S^3),\qquad
y_\infty=h_0p\in Z_\infty^{5,38}(S^0).
\end{gathered}
\]
The article applies the classical input and the two extensions
\[
d_2(h_5)=h_0h_4^2,\qquad
d_1^{f,E_2}(h_5)=h_2h_5,\qquad
d_2^{f,E_\infty}(h_0h_4^2)=h_0p
\]
with their specified no-crossing conditions to obtain
\[
d_3(h_2h_5)=h_0p.
\]
Its parameters are $n=r=2$, $m=1$, $l=2$, $e(f)=1$, $s=1$ and $t=35$.
\uses{paper:5.4,paper:4.14}
\lean{KIP126.ProofTools.Standard.paper6_7_classical_d3}
% Correspondence: The public theorem proves the exact nonzero d3 with the archived h0p target, using the registered independent classical input. It does not certify the original Leibniz derivation; its crossing comparisons are distinguished below.
\end{proposition}
\begin{proof}
\uses{sem:Literature:two-line-classical-inputs,sem:Computation:bounded-comparison}
\textit{An independent proof of the numerical conclusion.} The classical
input in \ref{sem:Literature:two-line-classical-inputs} gives a nonzero
$d_3$ from the specified $h_2h_5$ to degree $(5,38)$. The bounded Ext
archive identifies the unique basis vector in that degree as $h_0p$,
where $p$ is the fixed archived generator in degree $(4,37)$.
Consequently every nonzero target in this one-dimensional
$\mathbb F_2$ space is $h_0p$, proving the displayed differential.
This proof uses neither the infinite extension nor the generalized
Leibniz theorem.

\textit{Source argument and scope of the revised tool.} The classical differential is the Hopf-invariant-one differential. The finite extension is the product with $h_2$; the infinite extension is the assertion of Example~\ref{paper:6.24}. The article's length formula gives $r+l-m=3$.

To use Theorem~\ref{sem:ProofTools:generalized-leibniz-at-level} instead of the original rule, at level zero one must verify actual top-or-left no-crossing in the quotient of length $1$ and actual no-crossing on the propagated infinite bottom edge. The classical checks of Example~\ref{paper:5.13} are not automatically those conditions. The independent proof of the numerical conclusion above does not close these finite-to-infinite and crossing comparisons in the original proposed derivation.

An application of the finite-layer surjectivity composition would be
a separate derivation: it must identify the first-quotient connecting
class with the specified $h_2h_5$, produce its two classical
differentials and verify every instance of (Surj) in the actual
$\nu$ triangle. Neither the existence of that general composition
statement nor the original no-crossing check verifies these
numerical conditions. This entry does not assert that this
alternative application has already been completed.
\end{proof}

\section{The finite-quotient core and classical composition}

\begin{theorem}[Finite-quotient Mahowald comparison at a specified level]
\label{sem:ProofTools:ml-finite-core}
\level{coarse}
A project derivation from the finite May construction in
Lemma~\ref{paper:6.11}, the compatible normalized finite rows, and finite
Bockstein comparison. The May input remains the external result in
\ref{paper:3.20}; it is not a separate level-indexed classical theorem.\par
Work with the mod-$2$ Adams spectral sequence of $2$-complete, connective,
finite-type spectra, fixed coherent synthetic and finite Bockstein comparisons,
and separated finite-quotient filtrations.
Fix one actual cofiber triangle
\[
 X\xrightarrow fY\xrightarrow gZ\xrightarrow h\Sigma X
\]
with its compatible normalized triangle, and write
$e_f+e_g+e_h=1$, with each exponent in $\{0,1\}$.
Absorb the classical suspensions into the three objects. The first-quotient
operations of the normalized arrows then have degrees
\[
 f_1:(e_f,e_f),\qquad g_1:(e_g,e_g),\qquad h_1:(e_h,e_h-1).
\]
When an exponent is zero this is the ordinary page map (with the displayed
desuspension for $h$). When it is one this is the actual normalized
first-quotient operation. These are not three freely chosen matrices.
All finite quotient rows, restriction maps, lambda inclusions, boundary
maps and graded comparisons below use these same arrows and coherent shifts.

Put $Q_nW=\nu W/\lambda^n$. Fix integers $S,T$, a level $c\geq0$ and
an integer $K\geq2$, and set $p=c+K$. Suppose
\[
 z\in Z_{K-1}^{S-e_h,T-e_h+1}(Z),\qquad
 u=h_1z\in E_2^{S,T}(X),\qquad
 v\in E_2^{S+e_f+K-1,T+e_f+K-1}(Y),
\]
and the actual classical differential is
\[
 d_Kz=g_1v\pmod{B_{K-1}(Z)}.
\]
The cycle qualification of its target is part of this page equation.
Define the actual homotopy groups
\[
 A=\pi_{T-S,T-c}Q_pX,\qquad
 H=\pi_{T-S+1,T-c+1-e_h}Q_pZ,\qquad
 B=\pi_{T-S,T-c+e_f}Q_pY.
\]
Write $\widehat f_p:A\to B$ and $h_p:H\to A$ for the induced normalized
maps, retaining their shifts and the fixed suspension comparison.
Let $\iota_pv\in B$ be the image of the specified first-quotient label
under the actual top-layer inclusion
$\Sigma^{0,-(p-1)}Q_1Y\to Q_pY$; its label is $\lambda^{p-1}v$.
Then there exists an actual representative $a_M\in F^SA$ such that
\[
 \operatorname{gr}^{S}(a_M)=[\lambda^cu],\qquad
 \widehat f_p(a_M)=\iota_pv.
\tag{M}
\]
In particular $u\in Z_K(X)$; no nonzero-label hypothesis is imposed.
If, in addition, the actual higher-filtration ambiguity satisfies
\[
 \widehat f_p(F^{S+1}A)=0,
\tag{A}
\]
then every specified $a\in F^SA$ with
$\operatorname{gr}^{S}(a)=[\lambda^cu]$ satisfies
\[
 \widehat f_p(a)=\iota_pv.
\tag{Cores}
\]
This is an equality in the length-$p$ quotient, with its original target.
There is no infinite-lift or target-permanence hypothesis.
The zero-level version is the specialization $c=0$, $p=K$.
The general level is obtained by the actual coefficient inclusion
$\Sigma^{0,-c}Q_K\to Q_{c+K}$, retaining both labels; it does not
assert a reverse implication obtained by cancelling $\lambda^c$.
\uses{paper:3.20,paper:6.11,sem:AdamsSynthetic:normalized-map,sem:AdamsSynthetic:lambda-quotients,sem:Literature:normalized-finite-quotient-row,sem:Literature:infinite-finite-window}
\lean{KIP126.ProofTools.Standard.finiteQuotientMahowaldLeibnizCore}
% Correspondence: The public theorem proves the complete core in the three exhaustive normalization coordinates of the same fixed-model triangle. It retains every finite level, the actual shift comparisons, the specified-label May representative, the original source cycle qualification, and the exact image of every same-label representative under (A). Its conclusion also identifies the signed coefficient target with the single highest-layer inclusion. The proof inherits the registered upstream construction and comparison obligations.
\end{theorem}
\begin{proof}
\uses{paper:3.8,paper:4.6,paper:6.11,sem:Literature:normalized-finite-quotient-row,sem:Literature:infinite-finite-window,sem:ProofTools:finite-strip,sem:ProofTools:sphere-label-detection}
First construct a lift of the \emph{specified} $z$ to $Q_{K-1}Z$ whose
adjacent boundary is exactly $g_1v$ in $Q_1Z$. Finite Bockstein comparison
first supplies a lift with boundary label $w$. The input page equation
says $w-g_1v\in B_{K-1}$. Each earlier $d_r$ boundary, $2\leq r<K$,
is the adjacent boundary of a $Q_{r-1}$ representative. Its image under
$\lambda^{K-r}:Q_{r-1}\to Q_{K-1}$ has zero first-quotient restriction
and the same adjacent boundary label. Adding these correction classes
therefore changes $w$ to $g_1v$ while retaining $z$ exactly.
Equivalently, the adjacent boundary of the kernel of first-quotient
restriction has label space $B_{K-1}$. For $K=2$ this correction space
is $B_1=0$.

First apply the finite May construction to this actual boundary square
with quotient lengths $K-1$ and $1$, in the same normalized triangle.
It constructs a single $a_0$ in $Q_KX$ whose restriction is the normalized
$h$ image of the chosen lift and whose normalized $f$ image is
$\iota_Kv$. The triangle convention may put a minus sign on the latter;
it disappears on the image of the mod-$2$ first-quotient label.

The finite-window comparison and naturality for the normalized $h$
map identify the first-quotient label of $a_0$ with $h_1z=u$.
Restriction is injective on the retained bottom grade, so finite
filtration reflection puts $a_0$ in $F^S$, identifies its specified
grade, and qualifies $u$ in $Z_K$. This includes the zero-label case.
Now apply the actual coefficient inclusion
$\Sigma^{0,-c}Q_KX\to Q_{c+K}X$ to $a_0$, and call its image $a_M$.
Naturality with the same normalized arrows and the finite label
comparison carry the two labels to $\lambda^cu$ and
$\lambda^{c+K-1}v$. The latter is the highest target grade.
Its actual representative equals the single top-layer inclusion
$\iota_pv$: their difference has filtration at least
$S+e_f+K$, where the target finite quotient is zero.
Thus this identification is an equality of actual homotopy classes,
not just an agreement of leading labels.
Hence $a_M$ has precisely the degree, leading term
and image asserted in (M). For $c=0$ this operation is the identity.
This also explains why a positive $e_h$ requires normalized graded
naturality, rather than an assertion about an ordinary classical page map.
A positive $e_g$ acts here on $Q_1Y$; it introduces no additional survival
requirement on $v$.

Finally, two representatives of the same associated grade differ by an
element of $F^{S+1}A$. Apply (A) to $a-a_M$ and use additivity:
$\widehat f_p(a)-\widehat f_p(a_M)=0$. Equation (M) gives the specified
image in (Cores), without changing either representative or cancelling
any power of $\lambda$.
\end{proof}

\begin{lemma}[Two finite criteria for annihilating the ambiguity]
\label{sem:ProofTools:ml-ambiguity-criteria}
\level{coarse}
A finite-filtration consequence of \ref{sem:ProofTools:finite-correction}.\par
In the notation of \ref{sem:ProofTools:ml-finite-core}, the finite-window
comparison identifies, for $0\leq j<K$,
\[
 \operatorname{gr}^{S+j}A
 =Z_{K-j}^{S+j,T+j}(X)/B_{c+j+1},\qquad F^{S+K}A=0.
\]
The highest target grade is $S+e_f+K-1$, and
$F^{S+e_f+K}B=0$.
Either of the following is sufficient for (A).

\emph{Surjective graded lifts.} For each $1\leq j<K$ and
$a\in F^{S+j}A$, there exists
$b\in F^{S+j-e_h}H$ such that
\[
 a-h_p(b)\in F^{S+j+1}A.
\]

\emph{Collapsed intermediate grades.} For $1\leq j<K-1$ the grade
$\operatorname{gr}^{S+j}A$ is zero, and the map induced by $\widehat f_p$
from the entire last grade $\operatorname{gr}^{S+K-1}A$ to
$\operatorname{gr}^{S+e_f+K-1}B$ is zero.
The latter criterion does not require the last source grade to vanish
or to come from $h_p$.
\uses{sem:ProofTools:ml-finite-core,sem:ProofTools:finite-strip,sem:Literature:infinite-finite-window}
\lean{KIP126.ProofTools.Standard.finiteSurjectiveAmbiguityZero,KIP126.ProofTools.Standard.finiteCollapsedAmbiguityZero}
% Correspondence: The public theorems prove that the entire higher-filtration ambiguity has zero image under the two displayed criteria. The collapsed-grade theorem retains the target's terminal zero filtration and the image in that filtration, rather than asking its caller to supply actual image vanishing. The actual graded comparisons in the applications remain those specified in the statement.
\end{lemma}
\begin{proof}
\uses{sem:ProofTools:finite-correction}
For the first criterion, start with $a_1\in F^{S+1}A$. Choose successive
corrections $b_j$ so that $a_{j+1}=a_j-h_p(b_j)\in F^{S+j+1}A$.
After $K-1$ steps $a_K=0$, hence
$a_1=h_p(\sum_{j=1}^{K-1}b_j)$. The actual triangle gives
$\widehat f_ph_p=0$, so $\widehat f_p(a_1)=0$.
For the second criterion, the intermediate zero grades give
$F^{S+1}A=F^{S+K-1}A$. The zero map on the last grade sends this subgroup
into $F^{S+e_f+K}B$, which is zero by the finite strip and separation.
For $K=2$ there are no intermediate grades; the last-grade condition
alone is still required.
\end{proof}

\begin{theorem}[Mahowald--Leibniz composition under finite-layer surjectivity]
\label{sem:ProofTools:cprime}
\level{coarse}
A project derivation using the finite Mahowald construction, the finite
Bockstein comparison of \cite[Theorem A.1 and Corollary A.11]{BHS},
and the finite correction argument. This is not a numbered statement
of \cite{LWX126}.\par
Work with the mod-$2$ Adams spectral sequence of $2$-complete, connective,
finite-type spectra, with the usual coherent synthetic comparisons and
separated finite-quotient filtrations. Fix one actual cofiber triangle
\[
X\xrightarrow fY\xrightarrow gZ\xrightarrow h\Sigma X,
\qquad e_f=1,\quad e_g=e_h=0.
\]
Here $e_f=1$ means positive Adams filtration, not necessarily filtration
exactly one. Use the compatible normalized synthetic lift of this same
triangle. Its homology short exact sequence
\[
0\longrightarrow H_*Y\xrightarrow{g_*}H_*Z
 \xrightarrow{h_*}\Sigma H_*X\longrightarrow0
\]
defines the Ext connecting homomorphism
\[
\partial_{g,h}:E_2^{s,t}(X)\longrightarrow E_2^{s+1,t+1}(Y).
\]
Write $g_2$ for the ordinary $E_2$ map of $g$, and $h_2$ for the ordinary
$E_2$ map of $h$ followed by desuspension, so that $h_2$ lowers internal
degree by one. The operation $\partial_{g,h}$ is the first-quotient
operation of the compatible normalized $\widehat f$, not the ordinary
$E_2(f)$, which is zero. The actual short exact sequence is essential:
two abstract matrices $g_2,h_2$ alone do not specify $\partial_{g,h}$.

Use $Z_r,B_r\subseteq E_2$ for cycles and boundaries through $d_r$,
with $Z_1=E_2$ and $B_1=0$; all boundary labels are allowed.
Let $R,K\geq2$ be integers, $s,t\in\mathbb Z$, and put
\[
S=s+R,\qquad T=t+R-1,\qquad D=R+K-1.
\]
Suppose
\[
x\in Z_{R-1}^{s,t}(X),\qquad
z\in Z_{K-1}^{S,T+1}(Z),\qquad
v\in E_2^{S+K,T+K}(Y),
\]
with the following actual page equations, including qualification of
their target labels on the indicated pages:
\[
d_Rx=h_2z\pmod{B_{R-1}^{S,T}(X)},\qquad
d_Kz=g_2v\pmod{B_{K-1}^{S+K,T+K}(Z)}.
\]
For every $1\leq j\leq K-1$, assume the finite-layer surjectivity condition
\[
Z_{K-j}^{S+j,T+j}(X)
=h_2\bigl(Z_{K-j}^{S+j,T+j+1}(Z)\bigr)
 +B_{R+j-1}^{S+j,T+j}(X).
\tag{Surj}
\]
Then the specified class $y=\partial_{g,h}(x)$ satisfies
\[
y\in Z_{D-1}^{s+1,t+1}(Y),\qquad
d_Dy=v\pmod{B_{D-1}^{S+K,T+K}(Y)}.
\]
No $f$-extension, nonzero $y$, survival of $y$, or permanence of $v$
is an input. In particular, if $y=0$, then $v\in B_{D-1}$.
If $v\notin B_{D-1}$, the resulting differential and the source class
on $E_D$ are nonzero.

\textit{Relation with the level-indexed rules.}
The output is a classical Adams differential; the intermediate lambda
level in the proof is fixed by $c=R-2$. The index $j$ in (Surj)
runs through a finite set of higher filtration grades, rather than
being another freely chosen lambda level of the output.
The theorem does not assume an $f$-extension or an infinite lift.
Its finite Mahowald instance has an $h$ edge of length zero, for which
the actual crossing triangle $s<u\leq v\leq s$ is empty.
The remaining ambiguity is eliminated by (Surj), not by an
unrestricted comparison between classical and synthetic crossings.
Thus this classical composition does not require first proving the
original infinite-extension application of Example~\ref{paper:6.24}.
The cited May input remains a general external lemma; only its
finite-quotient instance is used here.
\uses{sem:AdamsSynthetic:lambda-quotients,sem:AdamsSynthetic:normalized-map,sem:AdamsSynthetic:homology-ext-boundary,paper:3.20,sem:Literature:normalized-finite-quotient-row}
\lean{KIP126.ProofTools.Standard.finiteLayerSurjectivityMahowaldLeibnizExt}
% Correspondence: The complete local proof first derives the same rotated homology SES from positive Adams filtration, then applies the proved fixed-model composition with every surjectivity layer and the registered upstream first-quotient/Ext comparison. No homology exactness or comparison equality is added as a caller premise. Five inherited upstream responsibilities remain registered: the canonical comodule abelian structure, homology-comodule recognition, E2/Ext coordinates, exactness lift, and special-fiber comparison. The local completion does not assert that these foundations have been constructed.
\end{theorem}
\begin{proof}
\uses{sem:AdamsSynthetic:homology-exact,sem:AdamsSynthetic:filtration-comparison,sem:AdamsSynthetic:homology-ext-boundary,sem:Literature:first-quotient-ext-boundary,sem:ProofTools:ml-finite-core,sem:ProofTools:ml-ambiguity-criteria,paper:4.6,sem:Literature:infinite-finite-window,sem:ProofTools:finite-strip}
First obtain the displayed underlying graded homology short exact
sequence from the same triangle. Positive Adams filtration of $f$,
by the factorization comparison, implies $H_*f=0$.
Tensor the distinguished triangle with $H\mathbb F_2$ and rotate its
homotopy long exact sequence. In every integer degree, exactness
and the vanishing of $H_*f$ and $H_*(\Sigma f)$ show that $g_*$ is
injective, its image is the kernel of $h_*$, and $h_*$ is surjective.
The final target is $H_*(\Sigma X)$, identified with $\Sigma H_*X$
by the canonical tensor--suspension comparison; the arrows are the
actual $g$ and $h$ of the original triangle.

Set $c=R-2$ and $p=c+K=R+K-2$. The actual quotient boundary of the
first-quotient label of $x$, after undoing its displayed suspension, gives
\[
 a_\delta\in A=\pi_{t-s-1,t+1}Q_pX.
\]
The finite Bockstein comparison for the given $d_R$ puts it in $F^SA$
and gives the specified leading grade $[\lambda^ch_2z]$.
Theorem~\ref{sem:ProofTools:ml-finite-core}, applied to the second
classical differential, constructs $a_M$ with that same grade and
$\widehat f_pa_M=\iota_pv$.

For $1\leq j<K$, finite-window comparison identifies the actual graded
$h_p$ map with
\[
 h_2:Z_{K-j}^{S+j,T+j+1}(Z)/B_{c+j+1}
       \longrightarrow Z_{K-j}^{S+j,T+j}(X)/B_{c+j+1}.
\]
Since $c+j+1=R+j-1$, (Surj) is exactly its surjectivity.
Lift each graded preimage to an actual representative in the same
filtration of $H$. The first criterion of
\ref{sem:ProofTools:ml-ambiguity-criteria} gives (A), including the
$j=K-1$ layer with denominator $B_p$. Thus
$\widehat f_pa_\delta=\iota_pv$.

The coherent quotient-boundary square identifies the left side with
$\delta_Y(\partial^Q_fx)$. Its equality with $\iota_pv$ decodes as
follows. Reduction to $Q_{p-1}Y$ kills this top-layer image; exactness
therefore lifts $y=\partial^Q_fx$ to $Q_pY$, giving $y\in Z_p$.
For a lift with adjacent boundary label $v_0$, naturality gives
$\iota_pv_0=\iota_pv$. The kernel of the actual top-layer inclusion
on these labels is $B_p$, by the finite-window comparison and vanishing
above the top filtration. Hence $v_0-v\in B_p$, and the adjacent
Bockstein criterion gives $d_{p+1}y=v$ modulo $B_p$.
The compatible first-quotient/Ext comparison identifies
$\partial^Q_f$ with $\partial_{g,h}$ of the same triangle.
Precisely, Lemma~\ref{sem:AdamsSynthetic:homology-ext-boundary}
lifts the already proved degreewise short exactness to the same comodule
sequence, and Theorem~\ref{sem:Literature:first-quotient-ext-boundary}
identifies its genuine Ext boundary with the selected operation.
These are the fixed upstream interfaces; neither group short exactness
nor equality of kernels and images alone establishes the identification.
The output bidegree is $(S+K,T+K)$, as required.
No target permanence or prior survival of $y$ was used.
\end{proof}

\begin{theorem}[Classical composition with collapsed intermediate layers]
\label{sem:ProofTools:ml-classical-nc}
\level{coarse}
A project derivation from \ref{sem:ProofTools:ml-finite-core}, with the
classical conditions (NC1) and (NC2) in \cite[Section 2]{MLConditions1020}.
The normalized maps and their coherent finite-quotient meaning are made
explicit here.\par
Use the triangle, normalized exponents and first-quotient operations of
\ref{sem:ProofTools:ml-finite-core}, allowing all three positions of the
positive exponent. Let $R,K\geq2$, $x\in Z_{R-1}^{s,t}(X)$,
$S=s+R$, $T=t+R-1$ and $p=R+K-2$. Suppose
\[
 z\in Z_{K-1}^{S-e_h,T-e_h+1}(Z),\quad
 u=h_1z\in E_2^{S,T}(X),\quad
 v\in E_2^{S+e_f+K-1,T+e_f+K-1}(Y),
\]
where $u=h_1z$ is an equality of the specified $E_2$ labels, and
\[
 d_Rx=u\pmod{B_{R-1}(X)},\qquad
 d_Kz=g_1v\pmod{B_{K-1}(Z)}.
\]
Assume the whole-space conditions
\[
 Z_{K-j}^{S+j,T+j}(X)\subseteq B_{R+j-1}^{S+j,T+j}(X)
 \quad(1\leq j<K-1),
\tag{NC1}
\]
and
\[
 f_1\bigl(E_2^{S+K-1,T+K-1}(X)\bigr)
 \subseteq B_p^{S+e_f+K-1,T+e_f+K-1}(Y).
\tag{NC2}
\]
Then the specified $y=f_1x\in E_2^{s+e_f,t+e_f}(Y)$ satisfies
\[
 y\in Z_p(Y),\qquad d_{p+1}y=v\pmod{B_p(Y)}.
\]
For $K=2$, (NC1) is empty and (NC2) still concerns the entire final
$E_2$ source. For $e_f=1,e_g=e_h=0$, this is the branch directly
comparable with \ref{sem:ProofTools:cprime}. The other two branches use
a normalized first-quotient operation on $g$ or $h$, respectively;
ordinary page maps cannot be substituted for those operations.
\uses{sem:ProofTools:ml-finite-core}
\lean{KIP126.ProofTools.Standard.classicalNCComposition}
% Correspondence: The public theorem proves the complete fixed-model ClassicalNCCompositionStatement, covering all three positions of the positive exponent, every NC1 layer and the entire final NC2 condition. Its actual normalized first-quotient operations are calibrated to ordinary page maps only when their exponent is zero. It produces the stated cycle qualification and specified differential without a caller-supplied representative, output-survival assumption or finite-square conclusion. It does not assert a separate Ext-boundary identification and inherits registered upstream construction and comparison obligations.
\end{theorem}
\begin{proof}
\uses{sem:ProofTools:ml-finite-core,sem:ProofTools:ml-ambiguity-criteria,paper:4.6,sem:Literature:infinite-finite-window,sem:ProofTools:finite-strip}
Take $c=R-2$ and the same specified $a_\delta$ and constructed $a_M$
as in the proof of \ref{sem:ProofTools:cprime}, now retaining the general
normalized shifts. Their difference lies in $F^{S+1}A$.
The quotient formula has denominator $B_{c+j+1}=B_{R+j-1}$, so (NC1)
annihilates the intermediate associated grades and gives
\[
 F^{S+1}A=F^{S+K-1}A.
\]
At the last layer, the source label space is $E_2/B_p$.
Normalized graded naturality identifies the map on that layer with
$f_1$ modulo the same $B_p$ in the highest target grade
$S+e_f+K-1$. By (NC2) this map is zero. Its actual image is therefore
in $F^{S+e_f+K}B=0$. This is precisely the second criterion of
\ref{sem:ProofTools:ml-ambiguity-criteria}, giving (A).
The core yields $\widehat f_pa_\delta=\iota_pv$.
The coherent boundary square and the same finite Bockstein decoding
then give $y\in Z_p$ and $d_{p+1}y=v$ modulo $B_p$.
The argument retains all three arrows, including the normalized
$g$ or $h$ operation when its exponent is positive.
\end{proof}

\begin{proposition}[The last-layer distinction between the two criteria]
\label{sem:ProofTools:ml-last-layer}
\level{coarse}
A project comparison of \ref{sem:ProofTools:cprime} and
\ref{sem:ProofTools:ml-classical-nc}.\par
In their common $e_f=1,e_g=e_h=0$ branch, put
\[
 U=E_2^{S+K-1,T+K-1}(X),\quad B_X=B_p(X),\quad B_Y=B_p(Y)
\]
in the corresponding source and target bidegrees, and let $\partial=f_1$.
Exactness gives $\ker\partial=\operatorname{im}h_2$.
The last Surj condition is equivalent to
\[
 \partial U=\partial B_X,
\]
whereas (NC2) only says $\partial U\subseteq B_Y$.
Naturality gives $\partial B_X\subseteq B_Y$, so last-layer Surj implies
(NC2). Under (NC2), last-layer Surj is equivalent to lifting all the
relevant target boundaries through the source boundary group:
\[
 \operatorname{im}\partial\cap B_Y=\partial B_X.
\]
In general the difference is measured by
\[
 \frac{\partial^{-1}(B_Y)}{\operatorname{im}h_2+B_X}
 \cong\frac{\operatorname{im}\partial\cap B_Y}{\partial B_X}.
\]
For an intermediate layer, (NC1) implies the corresponding Surj
condition by taking the $h_2$ preimage to be zero. At the level of these
additive-group conditions, the reverse implication need not hold: a
nonboundary cycle can have a cycle-qualified $h_2$ preimage. This is a
logical comparison of the conditions, not a claimed counterexample in a
specified spectrum. These comparisons alone establish neither implication
between the complete packages. The reverse last-layer implication requires
the displayed boundary-lifting condition.
\uses{sem:ProofTools:ml-finite-core,sem:ProofTools:ml-ambiguity-criteria,sem:AdamsSynthetic:first-quotient-shifts,sem:Literature:normalized-finite-quotient-row}
\lean{KIP126.ProofTools.Standard.mahowaldLeibnizLastLayerComparison,KIP126.ProofTools.Standard.mahowaldLeibnizLastLayerSurjectivity_iff,KIP126.ProofTools.intermediateLayerSurjectivity_not_converse}
% Correspondence: The actual fixed-model comparison derives kernel/image exactness and preservation of B_p from the same normalized triangle, with no caller-supplied exactness or boundary-preservation premise. It proves the complete quotient obstruction isomorphism and the cycle-qualified intermediate implication. The last-layer equivalence also identifies the actual Z_1 and DesuspendsClass formulation of Surj with the displayed whole-group decomposition. The final additive-group example witnesses only the stated logical nonconverse, not a counterexample in a selected spectrum or to either complete composition theorem. The proof uses the fixed first-quotient operation and does not require the separate Ext identification.
\end{proposition}
\begin{proof}
\uses{sem:ProofTools:ml-finite-core,sem:AdamsSynthetic:first-quotient-shifts,sem:Literature:normalized-finite-quotient-row}
Exactness of the same normalized first-quotient triangle gives
$\ker\partial=\operatorname{im}h_2$ after the fixed suspension comparisons.
The actual top-layer inclusion vanishes on a label exactly when that
label lies in $B_p$. Naturality of that inclusion under the normalized
map, together with the positive weight-shift comparison, therefore gives
$\partial B_X\subseteq B_Y$. At the last layer $Z_1=E_2$, and the
actual desuspension relation identifies its Surj condition with the
whole-group decomposition used below.

If $a=h_2z+b$ with $b\in B_X$, then $\partial a=\partial b$.
Conversely, if $\partial a=\partial b$ for such a $b$, exactness gives
$a-b=h_2z$. This proves the first equivalence. The restriction of
$\partial$ to $\partial^{-1}(B_Y)$ is onto
$\operatorname{im}\partial\cap B_Y$; taking its quotient by the
preimage of $\partial B_X$ gives the displayed isomorphism.
The intermediate-layer statement follows directly from the specified
cycle-qualified formulation of Surj. For its additive-group nonconverse,
take a nonzero group with identity $h_2$, all elements as cycles and zero
boundary subgroup. Every element has a cycle-qualified preimage, while
the cycle subgroup is not contained in the boundary subgroup. This
example asserts no realization by a particular selected spectrum.
\end{proof}

\begin{proposition}[The fixed-model Mahowald--Leibniz assertion under finite-layer surjectivity]
\label{sem:ProofTools:cprime-model}
\level{coarse}
Fix a triangle in the selected classical object family of the Standard
model, with $e_f=1$ and $e_g=e_h=0$.
Define $\partial^{Q}_{f}$ by transporting the actual map
\[
Q_1(\Sigma^{0,1}\nu X)\xrightarrow{Q_1(\widehat f)}Q_1(\nu Y)
\]
through the fixed first-quotient comparisons; its degree is $(1,1)$.
Define $\mathrm{Surj}(T,R,K,s,t)$ by (Surj) in
Theorem~\ref{sem:ProofTools:cprime}, using the actual $E_2$ map of $h$
and the fixed tower desuspension relation.
Define $\mathrm{ML}_{\mathrm{surj}}(T)$ to be the following proposition:
if the selected normalized triangle is distinguished, then, for every
$R,K\geq2$, $s,t$, and $x,z,v$ satisfying the two input cycle
qualifications, the two actual page equations and $\mathrm{Surj}$ of
that theorem, the class $y=\partial^Q_f(x)$ lies in $Z_{R+K-2}$ and
satisfies $d_{R+K-1}y=v$ with exactly its displayed bidegree.
This definition imposes no nonvanishing condition on $y$.
The proposition $\mathrm{ML}_{\mathrm{surj}}(T)$ holds for the selected
triangle and operations.

There are two distinct proof obligations in this description.
First, prove the fixed-model proposition using the constructed
$\partial^Q_f$, the two classical input differentials and (Surj).
Second, identify
\[
\partial^Q_f=\partial_{g,h}
\]
for the same actual homology short exact sequence, so that its output
is the connecting class specified in the classical theorem.
The operation is already defined from the selected normalized map;
the comparison with the Ext connecting homomorphism is not supplied
by that definition.

The finite constructions and residual correction prove the fixed-model
proposition. The Ext comparison remains separate. The bindings below
specify the proposition and its operations and prove the displayed
assertion; the proof theorem is distinct from the three definitions.
For a concrete application one must also provide its
actual input differentials, compatible triangle and every layer of
(Surj); an asserted output alone is not evidence for these inputs.
\uses{sem:ProofTools:ml-finite-core,sem:ProofTools:ml-ambiguity-criteria,sem:AdamsSynthetic:first-quotient-shifts,sem:AdamsSynthetic:normalized-map}
\lean{KIP126.ProofTools.Standard.cofiberFirstQuotientBoundary,KIP126.ProofTools.Standard.MahowaldLeibnizFiniteLayerSurjectivity,KIP126.ProofTools.Standard.FiniteLayerSurjectivityMahowaldLeibnizStatement,KIP126.ProofTools.Standard.finiteLayerSurjectivityMahowaldLeibniz}
% Correspondence: The first three declarations define the actual operation, conditions and fixed-model proposition. The fourth is its complete proof. The same selected classical object family, normalized triangle compatibility, input differentials and every surjectivity layer are retained. This theorem does not identify its first-quotient output with a separately constructed Ext boundary.
\end{proposition}
\begin{proof}
\uses{sem:ProofTools:ml-finite-core,sem:ProofTools:ml-ambiguity-criteria,paper:4.6,sem:Literature:infinite-finite-window,sem:ProofTools:finite-strip}
The finite Bockstein comparison constructs the representative of the
specified first-quotient boundary from the first input differential.
The common finite core constructs the May representative from the
second input differential. Their leading labels agree, and the complete
finite-layer surjectivity condition annihilates their difference by the
first ambiguity criterion. Decoding the resulting exact image equation
gives both the claimed output cycle and its specified differential.
Every operation in this argument is the fixed first-quotient operation;
no Ext identification is used.
\end{proof}

\section{Adams differential transport and classical realization of finite-level extensions}

The differential-transport rules have the shape
\[
d_bx=u,\qquad x\longmapsto y,\qquad u\longmapsto v
\quad\Longrightarrow\quad d_{b+L-M}y=v,
\]
where $M$ and $L$ are image filtration shifts. The quotient-boundary
square of generalized Leibniz, including the cutoff of
Remark~\ref{paper:6.5}, supplies sufficient hypotheses for this
conclusion. The complementary Mahowald construction takes an $h$
preimage, a differential, and a $g$ preimage, with $r=n+l+m$.
Classical extensions additionally require full representatives and
the realization of Proposition~\ref{paper:5.14}.

All extension equations below refer to the actual maps, specified
infinite-page classes and complete shorter-image indeterminacies of
\ref{sem:ProofTools:level-finite-extension}. The differential rules
come first: the zero-level Leibniz specialization, the cutoff square,
and its classical Adams consequence. The finite Mahowald specialization
then supplies the finite-level relations used in the classical
realization problem. Its output must satisfy the stated source
permanence and no-crossing hypotheses before that passage can be made.
The finite-layer surjectivity composition of \ref{sem:ProofTools:cprime}
has its own finite proof and does not require this classical realization.

\subsection{Adams differential transport}

\begin{corollary}[The zero-input-level Leibniz rule with actual crossing conditions]
\label{sem:ProofTools:ss-leibniz-specialization}
\level{coarse}
A specialization of
\ref{sem:ProofTools:generalized-leibniz-at-level} at input level zero.

In the fixed Standard model let $f:X\to Y$, $e=e(f)$ and
$F:\Sigma^{0,e}\nu X\to\nu Y$ be its actual normalized map. Suppose
\[
2\le n\le b,\qquad e\le M\le n-2+e,\qquad L\ge e.
\]
Take the specified labels
\[
\begin{gathered}
x\in E_2^{s,t}(X),\qquad y\in E_2^{s+M,t+M}(Y),\\
u\in Z_\infty^{s+b,t+b-1}(X),\qquad
v\in Z_\infty^{s+b+L,t+b+L-1}(Y).
\end{gathered}
\]
Assume the actual equations
\[
d_bx=u,\qquad \mathcal E^f_{n-1,0}(M;x,y),\qquad
d_L^F(u)=\lambda^{L-e}v.
\]
Put $p=(t-s,t+e)$ and $A=\Sigma^{0,e}\nu X$. Require
\[
\bigl(NC(F_{n-1};p,s,M)\lor
 NC(\delta_{n-1}^A;p,s,b)\bigr)
\land NC(\Sigma^{1,-(n-1)}F;p,s+b,L).
\]
Then
\[
d_{b+L-M}(y)=v.
\]
An original page-extension input $d_M^{f,E_n}(x)=y$ supplies the
middle equation by \ref{sem:ProofTools:extension-level-zero}; the
crossing conditions in this corollary remain conditions on the actual
maps at the displayed weight.
\uses{sem:ProofTools:generalized-leibniz-at-level,sem:ProofTools:extension-level-zero}
\lean{KIP126.ProofTools.Standard.ssLeibnizAtLevelZero}
% Correspondence: The public specialization retains all premises, including both permanent-label qualifications and the full bottom extension.
\end{corollary}
\begin{proof}
\uses{sem:ProofTools:generalized-leibniz-at-level,sem:ProofTools:extension-level-zero}
Set the lambda level to $0$, and substitute $r=b$, $m=M$ and $l=L$
in \ref{sem:ProofTools:generalized-leibniz-at-level}. Its quotient
exponent is $n-1$ and its weight is $t+e$, exactly as above. All its
equations and actual no-crossing premises are the displayed ones.
Its conclusion is the stated classical differential. This substitution
does not use the original crossing equivalence of \cite{LWX126}, Proposition 5.10.
\end{proof}

\begin{proposition}[The cutoff Leibniz square at an arbitrary lambda level]
\label{sem:ProofTools:ss-cutoff-square-specialization}
\level{coarse}
A quotient-boundary specialization of the four-edge theorem
\ref{paper:2.12} and Remark~\ref{paper:6.5}.

Fix $F:A\to B$ with $A=\Sigma^{0,e}\nu X$, $B=\nu Y$. For
$n\ge2$ and a lambda level $k\ge0$ put $q=n+k-1$. Use the actual
commutative square
\[
\begin{array}{ccc}
Q_qA&\xrightarrow{F_q}&Q_qB\\
\delta_q^A\downarrow&&\downarrow\delta_q^B\\
\Sigma^{1,-q}A&\xrightarrow{\Sigma^{1,-q}F}&\Sigma^{1,-q}B.
\end{array}
\]
Here $Q_qA=A/\lambda^q$. Fix a stem and weight $p=(a,w)$, and actual
infinite-page classes $\xi,\eta,\zeta,\omega$ in the four vertices,
at filtrations $s,s+M,s+b,s+b+L$, respectively; the remaining degree
coordinates are $(a+\mathrm{filtration},w)$.
Let $M,b,L\ge0$, $D=b+L-M$ and $1\le\kappa\le D$.
Assume
\[
d_M^{F_q}(\xi)=\eta,\qquad
d_b^{\delta_q^A}(\xi)=\zeta,\qquad
d_L^{\Sigma^{1,-q}F}(\zeta)=\omega.
\]
Require no crossing on the top or left equation. Require that the
bottom equation have no crossing hitting filtration at least
$s+M+\kappa$. Finally require the actual right-edge equation
\[
d_{\kappa-1}^{\delta_q^B}(\eta)=0
\]
with no crossing. Every crossing here is measured in the extension
spectral sequence of its actual edge, at the fixed stem and weight,
with its full shorter-image indeterminacy. Then
\[
d_D^{\delta_q^B}(\eta)=\omega.
\]
This is an extension-square conclusion. To read it as a specified
classical Adams differential, its labels must also have the required
boundary-comparison form. For example, use the classical labels,
degrees and range bounds of \ref{sem:ProofTools:ss-leibniz-specialization}.
If the actual comparisons give
\[
\begin{gathered}
\xi=\lambda^k x,\quad
\eta=\lambda^{k+M-e}y,\\
\zeta=\lambda^{b-n}u,\quad
\omega=\lambda^{b-n+L-e}v,
\end{gathered}
\]
at weight $t+e-k$, then the exponent identity
\[
D-q-1+(k+M-e)=b-n+L-e
\]
is precisely the arithmetic required for the right boundary to decode
as $d_Dy=v$. Those range bounds give $D\ge2$ and
$0\le k+M-e<q$. Actual qualified labels and comparisons remain
required; the exponent identity alone is not the decoding theorem.
\uses{paper:2.12,paper:6.5,sem:ProofTools:level-finite-extension}
% Blueprint note: The cutoff kappa and the lambda level k are different parameters. A zero differential range alone does not supply the crossing-free right-edge extension.
\lean{KIP126.ProofTools.Standard.ssCutoffSquareAtLevel}
% Correspondence: The public theorem proves the actual quotient-boundary conclusion at every level with its complete mathematical hypotheses. The exponent and range calculations are private auxiliary steps, not extra conclusion bindings.
\end{proposition}
\begin{proof}
\uses{paper:2.12}
Apply the four-edge theorem to this square, replacing its top length
by $M$, its left length by $b$, its bottom length by $L$, and its
cutoff by $\kappa$. The hypotheses include its right-edge zero
extension and both required crossing conditions. It returns length
$b+L-M=D$ on the right. The theorem applies to every actual quotient
$q=n+k-1$; setting $k=0$ gives the zero-input-level square with the
same cutoff hypotheses. The displayed exponent identity follows by
substitution of $D$ and $q$, while the actual boundary and label
comparisons remain distinct mathematical requirements.
\end{proof}

\begin{corollary}[Adams differential transport with a crossing cutoff]
\label{sem:ProofTools:ss-adams-soundness}
\level{coarse}
A classical Adams consequence of
\ref{sem:ProofTools:ss-cutoff-square-specialization}, using the
weighted boundary comparisons in the proof of
\ref{sem:ProofTools:generalized-leibniz-at-level}.

Work in the fixed Standard model. Let $f:X\to Y$, $e=e(f)$,
$A=\Sigma^{0,e}\nu X$, $B=\nu Y$, and let $F:A\to B$ be the
actual normalized map. Take
\[
2\le n\le b,\qquad e\le M\le n-2+e,\qquad L\ge e,\qquad k\ge0,
\]
and put $q=n+k-1$, $p=(t-s,t+e-k)$ and $D=b+L-M$.
Specify
\[
\begin{gathered}
x\in E_2^{s,t}(X),\qquad y\in E_2^{s+M,t+M}(Y),\\
u\in Z_\infty^{s+b,t+b-1}(X),\qquad
v\in Z_\infty^{s+b+L,t+b+L-1}(Y),
\end{gathered}
\]
with the actual relations
\[
d_bx=u,\qquad
\mathcal E^f_{q,k}(M;x,y),\qquad
d_L^F(u)=\lambda^{L-e}v.
\]
The middle relation has specified infinite-page source and target
classes in the finite quotients, with labels
$\lambda^kx$ and $\lambda^{k+M-e}y$. Denote its target class by
$\eta$. Boundary labels and zero labels are allowed.

Choose $1\le\kappa\le D$. On the actual quotient-boundary square
of \ref{sem:ProofTools:ss-cutoff-square-specialization}, require
\[
NC(F_q;p,s,M)\ \lor\ NC(\delta_q^A;p,s,b).
\]
For the propagated bottom map $\Sigma^{1,-q}F$, require that there
be no essential extension whose source and target filtrations
$a,c$ satisfy
\[
s+b<a\le c\le s+b+L,\qquad s+M+\kappa\le c.
\]
Finally require the right-edge zero extension from the same $\eta$,
\[
d_{\kappa-1}^{\delta_q^B}(\eta)=0,\qquad
NC(\delta_q^B;p,s+M,\kappa-1).
\]
All these conditions are at the fixed stem and weight $p$ in the
extension spectral sequences of the displayed actual maps. The
crossing ranges include length zero and their terminal target
filtrations. Then $D\ge2$ and the specified classical labels satisfy
\[
d_D(y)=v.
\]
This includes the cycle qualifications on $E_D$.
Setting $k=0$ gives the zero-input-level rule in quotient $Q_{n-1}$;
the cutoff $\kappa$ is independent of the lambda level $k$.
\uses{sem:ProofTools:level-finite-extension,sem:ProofTools:ss-cutoff-square-specialization}
\lean{KIP126.ProofTools.Standard.cutoffAdamsDifferentialAtLevel}
% Correspondence: The public theorem proves the complete classical Adams conclusion at every input level, including the actual right-edge zero extension from the specified target and the cutoff crossing ranges. Its representative and weighted boundary comparisons remain private proof steps.
\end{corollary}
\begin{proof}
\uses{sem:ProofTools:generalized-leibniz-at-level,sem:ProofTools:ss-cutoff-square-specialization,paper:4.6,sem:ProofTools:filtration-representatives}
Use the actual left-edge and bottom-edge constructions from the
proof of \ref{sem:ProofTools:generalized-leibniz-at-level}. The
differential $d_bx=u$ gives the left extension from the same source
class as the finite top extension. Its target is the propagated
class labelled by $\lambda^{b-n}u$. Multiplication by
$\lambda^{b-n}$ and the actual boundary suspension transport the
given infinite extension to the bottom edge, with specified target
label $\lambda^{b-n+L-e}v$. These constructions identify the actual
source and target classes, including their suspension comparisons.

The cutoff-square theorem now gives a length-$D$ right extension
from $\eta$ to this specified bottom target. Choose a representative
$a$ witnessing this extension. It detects $\lambda^d y$ in $Q_qB$,
where $d=k+M-e$, and its own actual boundary has the specified
bottom target detection. Undoing the boundary suspension gives a
representative in $\nu Y$ with the corresponding weighted label
of $v$.

For completeness, the bounds needed by the weighted reverse
boundary comparison are
\[
\begin{gathered}
0\le d<q,\qquad
q-d=n-1-M+e\le D-1,\qquad D\ge2,\\
(t+M)-d=t+e-k,\qquad
D-q-1+d=b-n+L-e.
\end{gathered}
\]
Apply that comparison to this same $a$ and its boundary, with quotient
length $q$, input degree $(s+M,t+M)$, input lambda exponent $d$ and
differential length $D$. It gives the actual Adams relation
$d_Dy=v$, with the required cycle qualifications. The displayed
identities verify the range and weight of these detected
representatives; they are used after the representative comparisons,
not as substitutes for them.
\end{proof}

\subsection{Finite-level extensions and classical realization}

\begin{corollary}[The finite Mahowald rule with both input levels zero]
\label{sem:ProofTools:ss-mahowald-specialization}
\level{coarse}
A specialization of the finite-layer revision
\ref{paper:6.12} at the two input levels zero.

Fix the same compatible normalized triangle and comparisons as in
\ref{paper:6.12}. Let
\[
\begin{gathered}
n_1=n-e_f\ge1,\quad m_1=m-e_g\ge0,\quad
k=l-e_h\ge0,\\
r=n+m+l,\qquad A=n_1+k,\quad B=m_1+1,\quad
r'=n_1+k+1.
\end{gathered}
\]
Thus $A+B=r$, $e_f+e_g+e_h=1$ and $r-k=n+m+e_h$.
Take the prescribed labels
\[
\begin{gathered}
\bar x\in Z_{r-1}^{s,t}(Z),\quad
\bar y\in Z_\infty^{s+r,t+r-1}(Z),\\
x\in Z_{n_1}^{s+l,t+l-1}(X),\quad
y\in Z_{m_1+1}^{s+l+n,t+l+n-1}(Y),
\end{gathered}
\]
with $d_r\bar x=\bar y$. For the actual finite maps and transported
labels of \ref{paper:6.12}, suppose
\[
d_l^{h_A}(\bar x)=\lambda^kx,\qquad
d_r^{\delta_{A,B}}(\bar x)=\lambda^{B-1}\bar y,\qquad
d_m^{g_B}(y)=\lambda^{B-1}\bar y.
\]
Require actual no-crossing on the $h_A$ equation or the
$\delta_{A,B}$ equation in the same weight. Their source class is
the same specified $\bar x$, and their common boundary target is
the same top-layer class, not just independently chosen labels.
Then there are the common actual representatives of
\ref{paper:6.12}, with
\[
\rho(c)=h_A(a),\qquad f_r(c)=-\lambda^Ab,
\]
and the conclusions
\[
x\in Z_{r-k},\qquad
d_n^{\widehat f_r}(\lambda^kx)=[\lambda^{r'-1}y].
\]
The source denominator is $B_{k+1}$ and the target denominator is
$B_{r'}$, with the shorter images of this same actual map also
retained. If $k=0$, this gives the ordinary level-zero page relation
in quotient $Q_r$, namely $\mathcal E^f_{r,0}(n;x,y)$, of effective
classical page $r+1$. If $k>0$, it retains the output level $k$;
the two zero input levels do not make its output level zero.
\uses{paper:6.12,sem:ProofTools:level-finite-extension}
\lean{KIP126.ProofTools.Standard.ssMahowaldZeroInputLevels}
% Correspondence: The public specialization includes the common representatives, signed equations, cycle improvement and retained output level. A smaller quotient or a full classical extension still requires its own division or lifting comparison.
\end{corollary}
\begin{proof}
\uses{paper:6.12}
Set $i=j=0$ in \ref{paper:6.12}. Its quotient lengths become
$A=n_1+k$ and $B=m_1+1$, its total length is $N=r$, its source
level is $K=k$, and its target level is $J=A=r'-1$. Its three
input equations, cycle qualifications, comparison choices and
no-crossing alternative are exactly those above. Hence its common
representatives, signed equations and specified output labels give
the asserted relation and $x\in Z_{r-k}$. When $k=0$ the source
and target windows are precisely those of the zero-level relation
in quotient $Q_r$. No division by $\lambda^k$ is used.
\end{proof}

\begin{theorem}[Permanence from a full lift of the specified finite target]
\label{sem:ProofTools:ss-full-lift-permanence}
\level{coarse}
A consequence of the full-to-finite comparison and filtration
reflection in the fixed model.

Let $N>0$, $y\in E_2^{s,t}(Y)$ and $0\le t-w<N$. Suppose an actual
full class
\[
\beta\in\pi_{t-s,w}(\nu Y)
\]
has quotient projection $q_N\beta$ detected by the specified finite
label $\lambda^{t-w}y$ in $Q_NY$. Then the original label satisfies
\[
y\in Z_\infty^{s,t}(Y).
\]
Neither nonvanishing nor prior permanence of $y$ is assumed. For a
classical $f$-extension, one additionally needs a single full source
representative whose actual normalized image has this prescribed
target detection, together with the realization comparisons of
Proposition~\ref{paper:5.14}.
\uses{sem:ProofTools:finite-strip,sem:ProofTools:filtration-representatives,sem:Literature:infinite-finite-window}
\lean{KIP126.ProofTools.finite_target_from_full_lift_is_permanent}
% Correspondence: The existing public theorem proves exactly the visible finite-label permanence assertion. It does not construct beta or identify an unrelated full source whose image is beta.
\end{theorem}
\begin{proof}
The prescribed finite detection puts $q_N\beta$ in filtration at
least $s$. Reflection along the full-to-finite projection in this
visible window puts $\beta$ in the same filtration. Its associated
grade is a full infinite-page class. Naturality and injectivity of
the finite-window comparison identify its classical label with $y$
up to the corresponding boundary group. The full label and these
boundary labels lie in $Z_\infty$, so $y$ lies there as well.
This criterion uses the specified full class $\beta$. The next
proposition constructs such a class as the actual image of a full
source representative.
\end{proof}

\begin{proposition}[Classical realization and target permanence of finite-level extensions]
\label{sem:ProofTools:ss-classical-soundness}
\level{coarse}
A target-preserving lifting statement using the finite-window
comparison and the classical realization of Proposition~\ref{paper:5.14}.
Work in the fixed Standard model, with the realization coordinates
and lambda-inversion Adams comparison of that proposition.

Let $f:X\to Y$, $e=e(f)$, $A=\Sigma^{0,e}\nu X$, $B=\nu Y$,
and let $F:A\to B$ be the actual normalized map. Take $R\ge2$,
$k\in\mathbb N$, and put
\[
q=R+k-1,\qquad p=(t-s,t+e-k).
\]
Suppose
\[
e\le m\le R-2+e,\qquad
x\in Z_\infty^{s,t}(X),\qquad
y\in E_2^{s+m,t+m}(Y),
\]
and assume the specified actual finite-level extension
\[
\mathcal E^f_{q,k}(m;x,y),\qquad
d_m^{F_q}(\lambda^kx)=\lambda^{k+m-e}y.
\]
Here the source is the intrinsic shifted finite label of
\ref{sem:ProofTools:level-finite-extension}; the equation includes
the indicated infinite-page classes and the complete shorter-image
indeterminacy of this same map.

Require $NC(F_q;p,s,m)$: there is no actual essential extension
$d_{v-u}^{F_q}(\xi)=\eta$ at this fixed stem and weight with
\[
s<u\le v\le s+m.
\]
This condition includes length zero and the terminal target
filtration $s+m$. Then
\[
y\in Z_\infty^{s+m,t+m}(Y),
\]
and there is one $a\in\pi_{t-s}(X)$ such that $x$ detects $a$ and
$y$ detects its actual image $f_*a$. Equivalently, the specified
permanent labels define the classical extension $d_m^f(x)=y$.
Neither essentiality of the given extension nor nonzero source
or target labels is required.

At $k=0$ this is a statement about a finite extension in
$Q_{R-1}$; an original $(f,E_R)$-extension supplies its finite
input by \ref{sem:ProofTools:extension-level-zero}. Positive
levels retain quotient length $R+k-1$ and weight $t+e-k$.
The proof retains the same full source representative through
the finite extension and classical realization at every level.
\uses{sem:ProofTools:level-finite-extension,sem:ProofTools:extension-level-zero,paper:5.14}
\lean{KIP126.ProofTools.Standard.finiteLevelExtensionToClassical}
% Correspondence: The public theorem proves target permanence and the detections of one classical source and its actual image at every level. The intrinsic source comparison uses only zero-level normalized-label calibration and invariance of realization under lambda multiplication. Existing upstream construction and comparison responsibilities remain inherited.
\end{proposition}
\begin{proof}
\uses{sem:ProofTools:extension-level-zero,sem:ProofTools:mh-actual-no-crossing,sem:ProofTools:ss-full-lift-permanence,sem:ProofTools:filtration-representatives,sem:Literature:infinite-nu-window,sem:Literature:infinite-finite-window,paper:5.14}
\textit{The full source at the prescribed level.}
Permanence of $x$ and the full $\nu$-window comparison give a
representative $\alpha_0\in\pi_{t-s,t}(\nu X)$ detected by $x$.
Its projection to $Q_1(\nu X)$ is the specified bottom label of
$x$. Apply the actual $(0,e)$ shift to obtain
$\widetilde\alpha_0\in\pi_{t-s,t+e}A$. The zero-level normalized
source comparison identifies its detection with the specified
shifted label of $x$. Hence the lambda-inversion Adams comparison
identifies its realization with a classical class $a$ detected
by the original $x$.

Put $\alpha=\lambda^k\widetilde\alpha_0$. Compatibility of
lambda multiplication with the actual shift identifies this
representative with the shift of $\lambda^k\alpha_0$. The native
lambda-label detection and the quotient-shift inclusion square
therefore show that $q_q^A\alpha$ is detected by the intrinsic
source label $\lambda^kx$ of the given finite extension.
Lambda multiplication preserves the realized homotopy class,
so this same $\alpha$ still realizes to $a$, detected by $x$.
Only the zero-level normalized-label calibration is used here;
the positive-level assertion follows from the actual lambda and
shift maps.

\textit{The actual image and its target detection.}
Choose the source filtration representative supplied by the
detection of $q_q^A\alpha$. Uniqueness of the infinite-page
representative identifies its grade with the source of the given
finite extension. Actual no crossing, including length zero and
the terminal target filtration, permits this particular source
representative while preserving the specified target grade.
Naturality of quotient projection then gives
\[
q_q^B(F_*\alpha)=F_{q*}(q_q^A\alpha)
\quad\text{detected by}\quad \lambda^{k+m-e}y.
\]
Set $\beta=F_*\alpha$. This constructs the required full target
lift from the same source. The bounds
\[
0\le k+m-e\le R+k-2=q-1
\]
place the target label in the visible finite window. By
\ref{sem:ProofTools:ss-full-lift-permanence}, $y$ is permanent;
filtration and detection reflection show that $\beta$ has the
specified full target detection at weight $t+e-k$.

\textit{Classical realization.}
The arbitrary-weight $\nu Y$ comparison identifies the realization
of $\beta$ as a class detected by $y$. The realization square for
the actual normalized map identifies this class with $f_*a$,
where $a$ is the realization of the same $\alpha$ above.
This proves both conclusions at every $k\ge0$, including $k=0$,
without a prior target permanence or target lifting hypothesis.
\end{proof}

The finite Mahowald theorem \ref{paper:6.12} constructs the finite
extension used in this proposition. In its notation the quotient
length and source level are $N$ and $K$, so the effective page
parameter here is $R=N-K+1$; its target level is
$J=K+n-e_f$. The source must additionally be permanent, and the
output map must satisfy the actual no-crossing condition above.
The no-crossing hypothesis on an input Mahowald edge does not by
itself assert this output condition. In particular, setting
$i=j=0$ in the Mahowald theorem can still leave the output at
level $k=l-e_h>0$. The general-level realization comparison is
then required. The finite-layer surjectivity composition
\ref{sem:ProofTools:cprime} instead proves its classical Adams
differential entirely in finite quotients.
